\chapter{Internal Categories and Mealy Morphisms in the Bicategory of Spans}
\label{ch:spans-internal-categories-mealy}

\section{Orientation}
\label{sec:ch1-orientation}

Let \(\mathcal E\) be a finitely complete category. Since \(\mathcal E\) has pullbacks, one may form the bicategory
\[
    \mathbf{Span}(\mathcal E),
\]
whose objects are those of \(\mathcal E\), whose 1-cells are spans, and whose 2-cells are maps of spans. In the formal theory of monads, internal categories in \(\mathcal E\) may be identified with monads in this bicategory. More explicitly, an endospan
\[
    C_0 \xleftarrow{s} C_1 \xrightarrow{t} C_0
\]
equipped with unit and multiplication as a monad in \(\mathbf{Span}(\mathcal E)\) is precisely an internal category in \(\mathcal E\) \cite{BenabouBicategories,StreetFormalTheoryOfMonads,LackCompanion,MirandaInternalCategories}.

The morphisms produced by this bicategorical identification require some care. We use the orientation in which a morphism from a monad \(A\) on \(A_0\) to a monad \(B\) on \(B_0\), carried by a 1-cell
\[
    P:A_0\to B_0,
\]
has structure 2-cell
\[
    \Phi:P\circ A\Rightarrow B\circ P.
\]
This orientation is commonly called a monad opmorphism, or a colax monad morphism, although the terminology varies between sources. We shall call it the \emph{Mealy orientation} and shall always display the structure cell explicitly.

In \(\mathbf{Span}(\mathcal E)\), the carrier \(P\) is a span
\[
    A_0 \xleftarrow{p} P \xrightarrow{q} B_0.
\]
By contrast, an internal functor
\[
    F:\mathbb A\to\mathbb B
\]
has object part a single morphism
\[
    F_0:A_0\to B_0.
\]
Thus the Mealy-oriented monad-morphism calculus and the ordinary functorial calculus for internal categories produce different kinds of morphism.

The present chapter takes the span-carried morphisms themselves as the object of study. We retain the span
\[
    A_0 \xleftarrow{p} P \xrightarrow{q} B_0
\]
as part of the structure and read it as carrying state. The map \(p\) records the input interface of a state, and \(q\) records its output interface. The oriented monad-morphism structure then supplies a transition-output operation. From this point of view, the span calculus suggests a typed internal form of automata theory: internal categories provide the input and output type systems, while monad morphisms in spans behave as stateful internal functors.

This interpretation is close in spirit to Par\'e's Mealy morphisms of enriched categories. Par\'e introduces Mealy morphisms as morphisms lying between enriched functors and profunctors, with the construction motivated by bicategorical considerations and by the classical Mealy-machine pattern from automata theory \cite{PareMealyMorphisms}. The present chapter develops an internal-categorical version of that idea in the bicategory of spans.

\section{Spans, monads, and internal categories}
\label{sec:spans-monads-internal-categories}

A span from \(X\) to \(Y\) is written
\[
    X \xleftarrow{p} S \xrightarrow{q} Y.
\]
The map \(p:S\to X\) is its left leg, and \(q:S\to Y\) is its right leg. We often display this as the diagram
\[
    \begin{tikzcd}
        & 
            S 
                \arrow[dl, "p"'] 
                \arrow[dr, "q"] 
        & 
        \\
            X 
        &
        & 
            Y.
    \end{tikzcd}
\]

When the structure maps are determined by context, we suppress them in pullback notation. For example, if
\[
    A_1\xrightarrow{t_A}A_0
    \qquad\text{and}\qquad
    P\xrightarrow{p}A_0
\]
are fixed, we write
\[
    A_1\times_{A_0}P
\]
for
\[
    A_1\times_{t_A,A_0,p}P.
\]
Thus \(A_1\times_{A_0}P\) is the pullback object in the square
\[
    \begin{tikzcd}
            A_1\times_{A_0}P
                \arrow[r, "\pi_P"]
                \arrow[d, "\pi_{A_1}"']
                \arrow[dr, phantom, "\lrcorner", very near start]
        &
            P
                \arrow[d, "p"]
        \\
            A_1
                \arrow[r, "t_A"']
        &
            A_0.
    \end{tikzcd}
\]
When several maps into the same object occur simultaneously, the relevant structure maps are displayed explicitly.

Let
\[
    X \xleftarrow{p} S \xrightarrow{q} Y
\]
and
\[
    Y \xleftarrow{r} T \xrightarrow{s} Z
\]
be spans in \(\mathcal E\). Their composite is obtained by forming the pullback of \(q:S\to Y\) and \(r:T\to Y\):
\[
    \begin{tikzcd}
            S\times_Y T
                \arrow[r, "\pi_T"]
                \arrow[d, "\pi_S"']
                \arrow[dr, phantom, "\lrcorner", very near start]
        &
            T
                \arrow[r, "s"]
                \arrow[d, "r"]
        & 
            Z
        \\
            S
                \arrow[r, "q"']
                \arrow[d, "p"']
        &
            Y 
        \\
            X.
    \end{tikzcd}
\]
We write this composite as \(T\circ S\), so \(S\) is performed first and \(T\)
second.

A map of spans
\[
    \alpha:
    \left(
        X 
        \xleftarrow{p} 
        S 
        \xrightarrow{q} 
        Y
    \right)
    \Longrightarrow
    \left(
        X 
        \xleftarrow{p'} 
        S' 
        \xrightarrow{q'} 
        Y
    \right)
\]
is a morphism \(\alpha:S\to S'\) in \(\mathcal E\) such that
\[
    p'\alpha=p,
    \qquad
    q'\alpha=q.
\]
Equivalently, it is a commuting diagram
\[
    \begin{tikzcd}[column sep=large]
        & 
            S 
                \arrow[dd, "\alpha" description] 
                \arrow[dl, "p"'] 
                \arrow[dr, "q"] 
        & 
        \\
            X 
        &
        & 
            Y 
        \\
        & 
            S' 
                \arrow[ul, "p'"] 
                \arrow[ur, "q'"'] 
        &
    \end{tikzcd}
\]
where the two triangular regions commute.

\begin{definition}[Bicategory of spans]
\label{def:bicategory-of-spans}
    The \textbf{bicategory of spans} in \(\mathcal E\), denoted
    \[
        \mathbf{Span}(\mathcal E),
    \]
    has
    \begin{enumerate}
        \item objects the objects of \(\mathcal E\);

        \item 1-cells \(X\to Y\) the spans
        \[
            X \xleftarrow{} S \xrightarrow{} Y;
        \]

        \item 2-cells the maps of spans.
    \end{enumerate}
    The identity 1-cell on \(X\) is
    \[
        X \xleftarrow{1_X} X \xrightarrow{1_X} X.
    \]
    Horizontal composition is given by the pullback construction above.
\end{definition}

\begin{proposition}
\label{prop:span-is-bicategory}
    If \(\mathcal E\) has pullbacks, then the data of Definition~\ref{def:bicategory-of-spans} form a bicategory.
\end{proposition}
\begin{proof}
    \leavevmode
    \begin{enumerate}
        \item \textbf{Composition of 2-cells.} Vertical composition is inherited from \(\mathcal E\): a vertical composite of maps of spans is the composite of their apex maps. The defining equations with the left and right legs are preserved by ordinary composition in \(\mathcal E\).

        Horizontal composition is induced by pullback. If
        \[
            \alpha:S\to S',
            \qquad
            \beta:T\to T'
        \]
        are maps of spans over a common middle object \(Y\), then the universal property of the pullback gives a unique map
        \[
            \alpha\times_Y\beta:S\times_YT\to S'\times_YT'
        \]
        whose projections are \(\alpha\pi_S\) and \(\beta\pi_T\). This is the horizontal composite of \(\alpha\) and \(\beta\).

        \item \textbf{Associators and unitors.}
        For composable spans with apexes \(S,T,U\), the two possible composite apexes
        \[
            (S\times_YT)\times_ZU
            \qquad\text{and}\qquad
            S\times_Y(T\times_ZU)
        \]
        both represent the object of triples \((s,t,u)\) satisfying the two matching equations. The associator is the canonical isomorphism between these two pullback presentations, uniquely determined by its projections to \(S,T,U\).

        The left and right unitors are the canonical pullback isomorphisms
        \[
            X\times_XS\cong S,
            \qquad
            S\times_YY\cong S.
        \]

        \item \textbf{Coherence.} The interchange law, the pentagon identity, and the triangle identity compare canonical maps between objects obtained from the same iterated pullback data. Since all such maps are uniquely determined by their projections to the original apexes, the relevant diagrams commute. Therefore the span data form a bicategory.
    \end{enumerate}
\end{proof}

\begin{lemma}[Component calculus for span maps]
\label{lem:component-calculus-span-maps}
    Let
    \[
        X\xleftarrow{}S\xrightarrow{}Y,
        \qquad
        X\xleftarrow{}T\xrightarrow{}Y
    \]
    be spans in \(\mathcal E\). A map of spans \(S\to T\) is determined by its apex map, and an equality of maps of spans is an equality of apex maps.

    Moreover, whenever the codomain apex is a pullback, a map into it is determined by its projections. Thus the generalized-element equations used below are equivalently morphism equations out of the pullback object that expresses the displayed typing constraints.
\end{lemma}

\begin{remark}
\label{rem:span-bicategory-literature}
    The construction of \(\mathbf{Span}(\mathcal E)\) is one of the basic examples in bicategory theory following B\'enabou's introduction of bicategories \cite{BenabouBicategories}.
\end{remark}

Let \(C_0\) be an object of \(\mathcal E\). The hom-category
\[
    \mathbf{Span}(\mathcal E)(C_0,C_0)
\]
is a monoidal category. Its objects are endospans
\[
    C_0\xleftarrow{s}S\xrightarrow{t}C_0,
\]
its morphisms are maps of spans, its tensor product is horizontal composition of spans, and its tensor unit is the identity span
\[
    C_0\xleftarrow{1_{C_0}}C_0\xrightarrow{1_{C_0}}C_0.
\]
We write this monoidal category as
\[
    \operatorname{EndSpan}_{\mathcal E}(C_0)
    :=
    \mathbf{Span}(\mathcal E)(C_0,C_0).
\]

With our convention for categorical order, if
\[
    C_0\xleftarrow{s}S\xrightarrow{t}C_0,
    \qquad
    C_0\xleftarrow{s'}T\xrightarrow{t'}C_0,
\]
then
\[
    S\otimes T:=T\circ S
\]
has apex
\[
    S\times_{t,C_0,s'}T.
\]
Thus \(S\otimes T\) denotes first \(S\), then \(T\):
\[
    \begin{tikzcd}
        & S \arrow[dl, "s"'] \arrow[dr, "t"] & & T \arrow[dl, "s'"'] \arrow[dr, "t'"] & \\
        C_0 & & C_0 & & C_0
    \end{tikzcd}
    \qquad\leadsto\qquad
    \begin{tikzcd}
        & S\times_{C_0}T \arrow[dl, "s\pi_S"']
                \arrow[dr, "t'\pi_T"] & \\
        C_0 && C_0
    \end{tikzcd}
\]

\begin{remark}
\label{rem:monads-as-monoids-in-endohom}
    In any bicategory \(\mathcal B\), a monad on an object \(X\) is equivalently a monoid object in the monoidal category
    \[
        \mathcal B(X,X),
    \]
    where the monoidal product is horizontal composition and the unit object is \(1_X\). Hence a monad on \(C_0\) in \(\mathbf{Span}(\mathcal E)\) is equivalently a monoid object in
    \[
        \operatorname{EndSpan}_{\mathcal E}(C_0).
    \]
    This is a basic instance of the formal theory of monads in bicategories \cite{StreetFormalTheoryOfMonads}.
\end{remark}

Let
\[
    C_0 \xleftarrow{s} C_1 \xrightarrow{t} C_0
\]
be an endospan. Composing this span with itself requires the pullback of
\[
    t:C_1\to C_0
    \qquad\text{and}\qquad
    s:C_1\to C_0.
\]
We denote this pullback by
\[
    C_2:=C_1\times_{t,C_0,s}C_1.
\]
Thus \(C_2\) is the object of composable pairs \((a,b)\) satisfying
\[
    t(a)=s(b).
\]
The defining pullback is
\[
    \begin{tikzcd}
            C_2
                \arrow[r, "\pi_2"]
                \arrow[d, "\pi_1"']
                \arrow[dr, phantom, "\lrcorner", very near start]
        &
            C_1
                \arrow[d, "s"]
        \\
            C_1
                \arrow[r, "t"']
        &
            C_0.
    \end{tikzcd}
\]
The composite endospan is
\[
    C_0
    \xleftarrow{s\pi_1}
    C_2
    \xrightarrow{t\pi_2}
    C_0.
\]

A monoid structure on the endospan
\[
    C_0 \xleftarrow{s} C_1 \xrightarrow{t} C_0
\]
consists of a unit map of spans
\[
    \eta:
    \left(
        C_0
        \xleftarrow{1_{C_0}}
        C_0
        \xrightarrow{1_{C_0}}
        C_0
    \right)
    \Longrightarrow
    \left(
        C_0
        \xleftarrow{s}
        C_1
        \xrightarrow{t}
        C_0
    \right)
\]
and a multiplication map of spans
\[
    \mu:
    \left(
        C_0
        \xleftarrow{s\pi_1}
        C_2
        \xrightarrow{t\pi_2}
        C_0
    \right)
    \Longrightarrow
    \left(
        C_0
        \xleftarrow{s}
        C_1
        \xrightarrow{t}
        C_0
    \right).
\]
Equivalently, the unit is a map
\[
    e:C_0\to C_1
\]
such that
\[
    se=1_{C_0},
    \qquad
    te=1_{C_0},
\]
as in
\[
    \begin{tikzcd}[column sep=large]
        & 
            C_0 
                \arrow[d, "e" description] 
                \arrow[dl, "1_{C_0}"'] 
                \arrow[dr, "1_{C_0}"] 
        & 
        \\
            C_0
        & 
            C_1 
                \arrow[l, "s"] 
                \arrow[r, "t"'] 
        & 
            C_0,
    \end{tikzcd}
\]
and the multiplication is a map
\[
    m:C_2\to C_1
\]
such that
\[
    sm=s\pi_1,
    \qquad
    tm=t\pi_2,
\]
as in
\[
    \begin{tikzcd}[column sep=large]
        & 
            C_2 
                \arrow[d, "m" description] 
                \arrow[dl, "s\pi_1"'] 
                \arrow[dr, "t\pi_2"] 
        & 
        \\
            C_0 
        & 
            C_1 
                \arrow[l, "s"] 
                \arrow[r, "t"'] 
        & 
            C_0.
    \end{tikzcd}
\]

In generalized-element notation, if
\[
    a:u\to v,
    \qquad
    b:v\to w,
\]
then
\[
    m(a,b)=b\circ a:u\to w.
\]
Thus the two projections from \(C_2\) are read as
\[
    \begin{tikzcd}[column sep=large]
            u
                \arrow[r, "a"]
        &
            v
                \arrow[r, "b"]
        &
            w.
    \end{tikzcd}
\]

\begin{definition}[Internal category]
\label{def:internal-category-as-endospan-monoid}
    An \textbf{internal category} in \(\mathcal E\) is a tuple
    \[
        \mathbb C=(C_0,C_1,s,t,e,m)
    \]
    where
    \[
        s,t:C_1\to C_0,
        \qquad
        e:C_0\to C_1,
        \qquad
        m:C_1\times_{t,C_0,s}C_1\to C_1,
    \]
    satisfying the source-target equations
    \[
        se=1_{C_0},
        \qquad
        te=1_{C_0},
    \]
    \[
        sm=s\pi_1,
        \qquad
        tm=t\pi_2,
    \]
    the identity equations
    \[
        m\circ\langle es,1_{C_1}\rangle
        =
        1_{C_1},
        \qquad
        m\circ\langle 1_{C_1},et\rangle
        =
        1_{C_1},
    \]
    and the associativity equation
    \[
        m\circ
        \left\langle
            m\circ\langle\pi_1,\pi_2\rangle,
            \pi_3
        \right\rangle
        =
        m\circ
        \left\langle
            \pi_1,
            m\circ\langle\pi_2,\pi_3\rangle
        \right\rangle
    \]
    as maps
    \[
        C_3\to C_1,
    \]
    where
    \[
        C_3
        :=
        C_1\times_{t,C_0,s}
        C_1\times_{t,C_0,s}
        C_1
    \]
    is the object of composable triples, with projections
    \[
        \pi_1,\pi_2,\pi_3:C_3\to C_1.
    \]
    All the displayed pairings are the unique maps into
    \[
        C_1\times_{t,C_0,s}C_1
    \]
    induced by its pullback universal property.
\end{definition}

\begin{theorem}[Endospan monoids are internal categories]
\label{thm:endospan-monoids-internal-categories}
    Let \(\mathcal E\) be a category with pullbacks. To give a monoid object in
    \[
        \operatorname{EndSpan}_{\mathcal E}(C_0)
    \]
    is equivalently to give an internal category in \(\mathcal E\) with object-of-objects \(C_0\).
\end{theorem}
\begin{proof}
    \leavevmode
    \begin{enumerate}
        \item \textbf{Unpack the monoid data.} A monoid object in \(\operatorname{EndSpan}_{\mathcal E}(C_0)\) is an
        endospan
        \[
            C_0\xleftarrow{s}C_1\xrightarrow{t}C_0
        \]
        together with a unit map of spans \(e:C_0\to C_1\) and a multiplication map of spans \(m:C_1\times_{t,C_0,s}C_1\to C_1\). The fact that these are maps of spans is exactly the source-target data
        \[
            se=1_{C_0},
            \qquad
            te=1_{C_0},
            \qquad
            sm=s\pi_1,
            \qquad
            tm=t\pi_2.
        \]

        \item \textbf{Identify the axioms.} The monoid unit axioms say that
        \[
            C_1
                \xrightarrow{\langle es,1\rangle}
            C_1\times_{C_0}C_1
                \xrightarrow{m}
            C_1,
            \qquad
            C_1
                \xrightarrow{\langle 1,et\rangle}
            C_1\times_{C_0}C_1
                \xrightarrow{m}
            C_1
        \]
        are both \(1_{C_1}\). These are precisely the internal identity laws. The monoid associativity axiom is the equality of maps \(C_3\to C_1\) which, on a composable triple
        \[
            \begin{tikzcd}[column sep=large]
                    u 
                        \arrow[r, "a"] 
                & 
                    v 
                        \arrow[r, "b"] 
                & 
                    w 
                        \arrow[r, "c"] 
                & 
                    z,
            \end{tikzcd}
        \]
        says
        \[
            c\circ(b\circ a)=(c\circ b)\circ a.
        \]
        Thus the monoid axioms are exactly the internal-category identity and associativity axioms.

        \item \textbf{Reverse the construction.} Conversely, an internal category \((C_0,C_1,s,t,e,m)\) determines the endospan
        \[
            C_0\xleftarrow{s}C_1\xrightarrow{t}C_0
        \]
        with unit \(e\) and multiplication \(m\). The internal-category axioms are exactly the monoid axioms just identified. Hence the two structures are equivalent.
    \end{enumerate}
\end{proof}

\begin{corollary}[Monads in spans are internal categories]
\label{cor:monads-in-spans-internal-categories}
    Let \(\mathcal E\) be a category with pullbacks. To give a monad in \(\mathbf{Span}(\mathcal E)\) on an object \(C_0\) is equivalently to give an internal category in \(\mathcal E\) with object-of-objects \(C_0\).
\end{corollary}
\begin{proof}
    \leavevmode
    \begin{enumerate}
        \item \textbf{Move from monads to monoids in an endohom.} By Remark~\ref{rem:monads-as-monoids-in-endohom}, monads on \(C_0\) in \(\mathbf{Span}(\mathcal E)\) are monoid objects in
        \[
            \operatorname{EndSpan}_{\mathcal E}(C_0)
            =
            \mathbf{Span}(\mathcal E)(C_0,C_0).
        \]

        \item \textbf{Apply the endospan-monoid identification.} By Theorem~\ref{thm:endospan-monoids-internal-categories}, monoid objects in this endohom category are exactly internal categories in \(\mathcal E\) whose object-of-objects is \(C_0\). This proves the claim.
    \end{enumerate}
\end{proof}

\begin{remark}
\label{rem:internal-categories-monads-literature}
    The identification of internal categories with monads in spans is a standard feature of the formal theory of internal categories \cite{StreetFormalTheoryOfMonads,LackCompanion,MirandaInternalCategories}. The present chapter uses this identification as the source of the stateful morphisms introduced below.
\end{remark}

\subsection{Two running specializations}
\label{subsec:two-running-specializations}

Two special cases will be carried through the chapter.

First, since \(\mathcal E\) is finitely complete, it has a terminal object \(1\) and finite products. With the standing choice of \(1\), binary products are realized as pullbacks over \(1\):
\[
    \begin{tikzcd}
            X\times Y
                \arrow[r]
                \arrow[d]
                \arrow[dr, phantom, "\lrcorner", very near start]
        &
            Y 
                \arrow[d, "!"]
        \\
            X 
                \arrow[r, "!"']
        &
            1.
    \end{tikzcd}
\]
A one-object internal category means an internal category whose object-of-objects is \(1\). Such an internal category is equivalently a monoid object in \(\mathcal E\).

We shall use the following order convention for this specialization. If
\[
    M=(M,e_M,\cdot_M)
\]
is a monoid object and multiplication is written by juxtaposition, then the associated one-object internal category
\[
    \mathbb M=(1,M,!,!,e_M,m_{\mathbb M})
\]
has internal-category composition
\[
    m_{\mathbb M}
    =
    \cdot_M\circ\tau_{M,M},
\]
where \(\tau_{M,M}:M\times M\to M\times M\) is the symmetry. Equivalently, for a composable pair \((a,b)\),
\[
    m_{\mathbb M}(a,b)=ba.
\]
Thus, in the one-object category,
\[
    b\circ a=ba.
\]
This convention is chosen so that the target-to-source execution law below becomes the usual right-action law
\[
    x\cdot(mn)=(x\cdot m)\cdot n.
\]

Thus the one-object specialization of the theory below is a theory of internal monoid-indexed Mealy systems.

Second, a morphism
\[
    F_0:A_0\to B_0
\]
determines a canonical span
\[
    \begin{tikzcd}
        & 
            A_0 
                \arrow[dl, "1_{A_0}"'] 
                \arrow[dr, "F_0"] 
        & 
        \\
            A_0 
        &
        & 
            B_0.
    \end{tikzcd}
\]
It also determines the reverse span
\[
    \begin{tikzcd}
        & 
            A_0 
                \arrow[dl, "F_0"'] 
                \arrow[dr, "1_{A_0}"] 
        & 
        \\
            B_0 
        &
        & 
            A_0.
    \end{tikzcd}
\]
The map-representable part of the theory recovers internal functors.
We denote these two spans by
\[
    F_{0!}
    :=
    \left(
        A_0
        \xleftarrow{1_{A_0}}
        A_0
        \xrightarrow{F_0}
        B_0
    \right)
\]
and
\[
    F_0^\ast
    :=
    \left(
        B_0
        \xleftarrow{F_0}
        A_0
        \xrightarrow{1_{A_0}}
        A_0
    \right),
\]
respectively.

\begin{definition}[Map-representable span]
\label{def:representable-span}
    A span
    \[
        A_0\xleftarrow{p}P\xrightarrow{q}B_0
    \]
    is called \textbf{map-representable}, or a \textbf{map span}, if it is isomorphic, as a span, to a span of the form
    \[
        A_0
        \xleftarrow{1_{A_0}}
        A_0
        \xrightarrow{F_0}
        B_0
    \]
    for some map \(F_0:A_0\to B_0\). Equivalently, the span is map-representable iff \(p\) is an isomorphism. In that case it is isomorphic to
    \[
        A_0\xleftarrow{1_{A_0}}A_0\xrightarrow{q p^{-1}}B_0.
    \]
    The isomorphism of spans is displayed by
    \[
        \begin{tikzcd}[column sep=large]
            & 
                P 
                    \arrow[d, "p" description, "\cong"'] 
                    \arrow[dl, "p"'] 
                    \arrow[dr, "q"] 
            & 
            \\
                A_0 
            & 
                A_0 
                    \arrow[l, "1_{A_0}"] 
                    \arrow[r, "qp^{-1}"'] 
            & 
                B_0.
        \end{tikzcd}
    \]
\end{definition}

\begin{proposition}[Canonical adjunction associated to a map]
\label{prop:canonical-span-adjunction}
    There is a canonical adjunction
    \[
        F_{0!}\dashv F_0^\ast
    \]
    in \(\mathbf{Span}(\mathcal E)\).
\end{proposition}
\begin{proof}
    \leavevmode
    \begin{enumerate}
        \item \textbf{Compute the composite \(F_0^\ast\circ F_{0!}\).} The composite
        \[
            F_0^\ast\circ F_{0!}:A_0\to A_0
        \]
        has apex the pullback of \(F_0\) with itself:
        \[
            \begin{tikzcd}
                    A_0\times_{B_0}A_0
                        \arrow[r, "\pi_2"]
                        \arrow[d, "\pi_1"']
                        \arrow[dr, phantom, "\lrcorner", very near start]
                &
                    A_0 
                        \arrow[d, "F_0"]
                \\
                    A_0 
                        \arrow[r, "F_0"']
                &
                    B_0.
            \end{tikzcd}
        \]
        Hence
        \[
            F_0^\ast\circ F_{0!}
            =
            \left(
                A_0
                \xleftarrow{\pi_1}
                A_0\times_{B_0}A_0
                \xrightarrow{\pi_2}
                A_0
            \right).
        \]

        \item \textbf{Define the unit.} The unit
        \[
            \eta_F:1_{A_0}\Rightarrow F_0^\ast\circ F_{0!}
        \]
        is induced by the diagonal map
        \[
            \Delta:A_0\to A_0\times_{B_0}A_0.
        \]
        The defining diagram is
        \[
            \begin{tikzcd}[column sep=large]
                & 
                    A_0 
                        \arrow[d, "\Delta" description] 
                        \arrow[dl, "1_{A_0}"'] 
                        \arrow[dr, "1_{A_0}"] 
                & 
                \\
                    A_0
                &
                    A_0\times_{B_0}A_0
                        \arrow[l, "\pi_1"]
                        \arrow[r, "\pi_2"']
                &
                    A_0.
            \end{tikzcd}
        \]

        \item \textbf{Compute the composite \(F_{0!}\circ F_0^\ast\).} The composite
        \[
            F_{0!}\circ F_0^\ast:B_0\to B_0
        \]
        has apex canonically isomorphic to \(A_0\), since the relevant pullback is
        \[
            \begin{tikzcd}
                    A_0\times_{A_0}A_0
                        \arrow[r]
                        \arrow[d]
                        \arrow[dr, phantom, "\lrcorner", very near start]
                &
                    A_0 
                        \arrow[d, "1_{A_0}"]
                \\
                    A_0 
                        \arrow[r, "1_{A_0}"']
                &
                    A_0.
            \end{tikzcd}
        \]
        Under this identification,
        \[
            F_{0!}\circ F_0^\ast
            =
            \left(
                B_0
                \xleftarrow{F_0}
                A_0
                \xrightarrow{F_0}
                B_0
            \right).
        \]

        \item \textbf{Define the counit.} The counit
        \[
            \varepsilon_F:F_{0!}\circ F_0^\ast\Rightarrow 1_{B_0}
        \]
        is induced by the map
        \[
            F_0:A_0\to B_0.
        \]
        It is displayed as
        \[
            \begin{tikzcd}[column sep=large]
                & 
                    A_0 
                        \arrow[d, "F_0" description] 
                        \arrow[dl, "F_0"'] 
                        \arrow[dr, "F_0"] 
                & 
                \\
                    B_0 
                & 
                    B_0 
                        \arrow[l, "1_{B_0}"] 
                        \arrow[r, "1_{B_0}"'] 
                & 
                    B_0.
            \end{tikzcd}
        \]

        \item \textbf{Verify the triangle identities.} The two triangle identities reduce to the equations defining the diagonal and the terminal pullback isomorphisms
        \[
            A_0\times_{A_0}A_0\cong A_0.
        \]
        Equivalently, both composites are canonical pullback-induced maps with the same projections. By uniqueness in the relevant pullback universal properties, the triangle identities hold.
    \end{enumerate}
\end{proof}

\section{Mealy morphisms}
\label{sec:mealy-morphisms}

Let
\[
    \mathbb A=(A_0,A_1,s_A,t_A,e_A,m_A)
\]
and
\[
    \mathbb B=(B_0,B_1,s_B,t_B,e_B,m_B)
\]
be internal categories in \(\mathcal E\), regarded as monads in \(\mathbf{Span}(\mathcal E)\):
\[
    \begin{tikzcd}
        & 
            A_1 
                \arrow[dl, "s_A"'] 
                \arrow[dr, "t_A"] 
        & 
        \\
            A_0 
        &
        & 
            A_0
    \end{tikzcd}
    \qquad
    \begin{tikzcd}
        & 
            B_1 
                \arrow[dl, "s_B"'] 
                \arrow[dr, "t_B"] 
        & 
        \\
            B_0 
        &
        & 
            B_0.
    \end{tikzcd}
\]

Let
\[
    A_0\xleftarrow{p}P\xrightarrow{q}B_0
\]
be a span:
\[
    \begin{tikzcd}
        & 
            P 
                \arrow[dl, "p"'] 
                \arrow[dr, "q"] 
        & 
        \\
            A_0 
        &
        & 
            B_0.
    \end{tikzcd}
\]
There are two composite spans associated with \(P\), \(A\), and \(B\).

First, the composite \(P\circ A\) is obtained by pulling back
\[
    t_A:A_1\to A_0
    \qquad
    \text{and}
    \qquad
    p:P\to A_0.
\]
Thus
\[
    \begin{tikzcd}
            A_1\times_{t_A,A_0,p}P
                \arrow[r, "\pi_P"]
                \arrow[d, "\pi_{A_1}"']
                \arrow[dr, phantom, "\lrcorner", very near start]
        &
            P 
                \arrow[d, "p"]
        \\
            A_1 
                \arrow[r, "t_A"']
        &
            A_0,
    \end{tikzcd}
\]
and
\[
    P\circ A
    =
    \left(
        A_0
        \xleftarrow{s_A\pi_{A_1}}
        A_1\times_{t_A,A_0,p}P
        \xrightarrow{q\pi_P}
        B_0
    \right).
\]

Second, the composite \(B\circ P\) is obtained by pulling back
\[
    q:P\to B_0
    \qquad
    \text{and}
    \qquad
    s_B:B_1\to B_0.
\]
Thus
\[
    \begin{tikzcd}
            P\times_{q,B_0,s_B}B_1
                \arrow[r, "\pi_{B_1}"]
                \arrow[d, "\pi_P"']
                \arrow[dr, phantom, "\lrcorner", very near start]
        &
            B_1 
                \arrow[d, "s_B"]
        \\
            P 
                \arrow[r, "q"']
        &
            B_0,
    \end{tikzcd}
\]
and
\[
    B\circ P
    =
    \left(
        A_0
        \xleftarrow{p\pi_P}
        P\times_{q,B_0,s_B}B_1
        \xrightarrow{t_B\pi_{B_1}}
        B_0
    \right).
\]

\begin{definition}[Mealy morphism]
\label{def:mealy-morphism}
    Let \(\mathbb A\) and \(\mathbb B\) be internal categories, regarded as monads in \(\mathbf{Span}(\mathcal E)\). A \textbf{Mealy morphism}
    \[
        (P,\Phi):\mathbb A\rightsquigarrow\mathbb B
    \]
    consists of a span
    \[
        A_0\xleftarrow{p}P\xrightarrow{q}B_0
    \]
    and a 2-cell in the Mealy orientation
    \[
        \Phi:P\circ A\Rightarrow B\circ P
    \]
    satisfying the unit and multiplication compatibility axioms
    \begin{align}
        \Phi\circ(P\eta_A)
        &=
        \eta_B P,
        \label{eq:mealy-unit-abstract}
        \\
        \Phi\circ(P\mu_A)
        &=
        (\mu_B P)\circ(B\Phi)\circ(\Phi A).
        \label{eq:mealy-multiplication-abstract}
    \end{align}
    The first equation is an equality of 2-cells
    \[
        P\Rightarrow B\circ P,
    \]
    and the second is an equality of 2-cells
    \[
        P\circ A\circ A\Rightarrow B\circ P.
    \]
    The equations are read with the canonical associativity and unit isomorphisms of \(\mathbf{Span}(\mathcal E)\) inserted where necessary.

    In terminology that calls a structure cell
    \[
        P\circ A\Rightarrow B\circ P
    \]
    a monad opmorphism or colax monad morphism, a Mealy morphism is precisely such a morphism carried by a span.
\end{definition}

The 2-cell
\[
    \Phi:P\circ A\Rightarrow B\circ P
\]
is a map
\[
    \Phi:
    A_1\times_{t_A,A_0,p}P
    \longrightarrow
    P\times_{q,B_0,s_B}B_1
\]
making the following span-map diagram commute:
\[
    \begin{tikzcd}[column sep=huge]  
        &
            A_1\times_{A_0}P
                \arrow[dl, "s_A\pi_{A_1}"']
                \arrow[dr, "q\pi_P"]
                \arrow[d, "\Phi" description]
        &  
        \\
            A_0
        &
            P\times_{B_0}B_1
                \arrow[l, "p\pi_P"]
                \arrow[r, "t_B\pi_{B_1}"']
        &
            B_0.
    \end{tikzcd}
\]
Write
\[
    \Phi(a,x)=\bigl(\delta(a,x),\omega(a,x)\bigr).
\]
Then the pullback condition says
\[
    s_B(\omega(a,x))=q(\delta(a,x)),
\]
and the two span-compatibility equations say
\[
    p(\delta(a,x))=s_A(a),
    \qquad
    t_B(\omega(a,x))=q(x).
\]
Equivalently, a Mealy morphism consists of maps
\[
    \delta:
    A_1\times_{t_A,A_0,p}P\to P,
    \qquad
    \omega:
    A_1\times_{t_A,A_0,p}P\to B_1,
\]
satisfying
\begin{align}
    p\delta&=s_A\pi_{A_1},
    \label{eq:mealy-definition-pdelta}
    \\
    s_B\omega&=q\delta,
    \label{eq:mealy-definition-source}
    \\
    t_B\omega&=q\pi_P.
    \label{eq:mealy-definition-target}
\end{align}
These equations are summarized by the diagram
\[
\begin{gathered}
\begin{tikzcd}[column sep=large, row sep=large]
        A_1\times_{A_0}P
            \arrow[r, "\delta"]
            \arrow[d, "\pi_{A_1}"']
    &
        P
            \arrow[d, "p"]
    \\
        A_1
            \arrow[r, "s_A"']
    &
        A_0
\end{tikzcd}
\qquad
\begin{tikzcd}[column sep=large, row sep=large]
        A_1\times_{A_0}P
            \arrow[r, "\omega"]
            \arrow[d, "\delta"']
    &
        B_1
            \arrow[d, "s_B"]
    \\
        P
            \arrow[r, "q"']
    &
        B_0
\end{tikzcd}
\qquad
\begin{tikzcd}[column sep=large, row sep=large]
        A_1\times_{A_0}P
            \arrow[r, "\omega"]
            \arrow[d, "\pi_P"']
    &
        B_1
            \arrow[d, "t_B"]
    \\
        P
            \arrow[r, "q"']
    &
        B_0.
\end{tikzcd}
\end{gathered}
\]
More transparently, an input arrow
\[
    \begin{tikzcd}[column sep=large]
            u 
                \arrow[r, "a"] 
        & 
            v
    \end{tikzcd}
\]
admissible at a state \(x\) with \(p(x)=v\) updates the state to one over \(u\), while producing an output arrow
\[
\begin{tikzcd}[column sep=huge]
        {q(\delta(a,x))}
            \arrow[r, "{\omega(a,x)}"]
    &
        {q(x)}.
\end{tikzcd}
\]
Thus the state update and the output arrow follow the target-to-source direction imposed by the structure cell
\[
    \Phi:P\circ A\Rightarrow B\circ P.
\]

\begin{proposition}[The Mealy laws]
\label{prop:mealy-laws}
    For a Mealy morphism
    \[
        (P,\delta,\omega):\mathbb A\rightsquigarrow\mathbb B,
    \]
    the abstract unit and multiplication axioms \eqref{eq:mealy-unit-abstract} and \eqref{eq:mealy-multiplication-abstract} are equivalent to the following component equations.

    The unit law is
    \begin{align}
        \delta(e_A(p x),x)&=x,
        \label{eq:delta-unit}
        \\
        \omega(e_A(p x),x)&=e_B(qx).
        \label{eq:omega-unit}
    \end{align}

    The multiplication law is
    \begin{align}
        \delta(b\circ a,x)
        &=
        \delta(a,\delta(b,x)),
        \label{eq:delta-composition}
        \\
        \omega(b\circ a,x)
        &=
        \omega(b,x)\circ \omega(a,\delta(b,x)),
        \label{eq:omega-composition}
    \end{align}
    for internally composable arrows
    \[
        a:u\to v,
        \qquad
        b:v\to w,
    \]
    and states \(x\in P\) satisfying \(p(x)=w\).
\end{proposition}
\begin{proof}
    \leavevmode
    \begin{enumerate}
        \item \textbf{Unpack the unit input.} The unit of \(\mathbb A\) induces a map
        \[
            P\to A_1\times_{A_0}P,
            \qquad
            x\mapsto (e_A(px),x).
        \]
        This is the map in the pullback diagram
        \[
            \begin{tikzcd}[column sep=large, row sep=large]
                P
                    \arrow[rr, "{\langle e_Ap,1_P\rangle}"]
                    \arrow[dr, "1_P"']
                    \arrow[ddr, "e_Ap"', bend right=15]
                &
                &
                A_1\times_{t_A,A_0,p}P
                    \arrow[dl, "\pi_P"]
                    \arrow[ddl, "\pi_{A_1}", bend left=15]
                \\
                &
                P
                    \arrow[d, "p"]
                \\
                &
                A_0
            \end{tikzcd}
        \]
        where \(t_Ae_Ap=p\).

        \item \textbf{Compare the two unit composites.} The unit axiom says that
        \[
            P
            \xrightarrow{(e_Ap,1_P)}
            A_1\times_{A_0}P
            \xrightarrow{\Phi}
            P\times_{B_0}B_1
        \]
        equals
        \[
            P
            \xrightarrow{(1_P,e_Bq)}
            P\times_{B_0}B_1.
        \]
        In diagram form:
        \[
            \begin{tikzcd}[column sep=large]
                    P
                        \arrow[r, "{(e_Ap,1_P)}"]
                        \arrow[dr, "{(1_P,e_Bq)}"']
                &
                    A_1\times_{A_0}P
                        \arrow[d, "\Phi"]
                \\
                &
                    P\times_{B_0}B_1.
            \end{tikzcd}
        \]
        Comparing the two components of the pullback \(P\times_{B_0}B_1\) gives
        \[
            \delta(e_A(px),x)=x,
            \qquad
            \omega(e_A(px),x)=e_B(qx).
        \]

        \item \textbf{Identify the multiplication domain.} For the multiplication axiom, the relevant generalized elements are triples
        \[
            (a,b,x)
        \]
        satisfying
        \[
            t_A(a)=s_A(b),
            \qquad
            t_A(b)=p(x).
        \]
        They are displayed as
        \[
            \begin{tikzcd}[column sep=large]
                    u 
                        \arrow[r, "a"] 
                & 
                    v 
                        \arrow[r, "b"] 
                & 
                    w
            \end{tikzcd}
            \qquad
            \text{with}\qquad
            p(x)=w.
        \]
        The corresponding pullback object is
        \[
            A_1\times_{t_A,A_0,s_A}A_1\times_{t_A,A_0,p}P.
        \]

        \item \textbf{Compute the left-hand side.} The left-hand side of \eqref{eq:mealy-multiplication-abstract} first composes \(a\) and \(b\) in \(\mathbb A\), then applies \(\Phi\):
        \[
            \begin{tikzcd}[column sep=large]
                    (a,b,x)
                        \arrow[r, maps to]
                &
                    (b\circ a,x)
                        \arrow[r, maps to]
                &
                    \bigl(\delta(b\circ a,x),\omega(b\circ a,x)\bigr).
            \end{tikzcd}
        \]

        \item \textbf{Compute the right-hand side.} The right-hand side first applies \(\Phi\) to \((b,x)\):
        \[
            (b,x)
            \longmapsto
            \bigl(\delta(b,x),\omega(b,x)\bigr).
        \]
        Since
        \[
            p(\delta(b,x))=s_A(b)=t_A(a),
        \]
        the pair \((a,\delta(b,x))\) is admissible. Applying \(\Phi\) again gives
        \[
            (a,\delta(b,x))
            \longmapsto
            \bigl(\delta(a,\delta(b,x)),\omega(a,\delta(b,x))\bigr).
        \]
        The two produced \(\mathbb B\)-arrows are composable:
        \[
            \begin{tikzcd}[column sep=huge]
                    {q(\delta(a,\delta(b,x)))}
                        \arrow[r, "{\omega(a,\delta(b,x))}"]
                &
                    {q(\delta(b,x))}
                        \arrow[r, "{\omega(b,x)}"]
                &
                    {q(x)}.
            \end{tikzcd}
        \]
        Hence their composite is
        \[
            \omega(b,x)\circ\omega(a,\delta(b,x)).
        \]

        \item \textbf{Compare components.} Equality of the two maps into \(P\times_{B_0}B_1\) is equality of the two \(P\)-components and the two \(B_1\)-components. Therefore the multiplication axiom is equivalent to
        \[
            \delta(b\circ a,x)
            =
            \delta(a,\delta(b,x)),
        \]
        and
        \[
            \omega(b\circ a,x)
            =
            \omega(b,x)\circ \omega(a,\delta(b,x)).
        \]
    \end{enumerate}
\end{proof}

Equivalently, using \(m_A(a,b)=b\circ a\), the multiplication laws may be written
\begin{align}
    \delta(m_A(a,b),x)
    &=
    \delta(a,\delta(b,x)),
    \label{eq:delta-composition-m}
    \\
    \omega(m_A(a,b),x)
    &=
    m_B\bigl(\omega(a,\delta(b,x)),\omega(b,x)\bigr).
    \label{eq:omega-composition-m}
\end{align}

\begin{remark}[Contravariant input processing]
\label{rem:contravariant-input-processing}
    The input is processed target-to-source. If
    \[
        a:u\to v
    \]
    and \(p(x)=v\), then
    \[
        p(\delta(a,x))=u.
    \]
    Operationally, a Mealy morphism may be regarded as a right-action-like contravariant input process on the state span, together with an output cocycle valued in the output category.
\end{remark}

\subsection{The one-object specialization}
\label{subsec:one-object-mealy-morphisms}

Let \(M\) and \(N\) be monoid objects in \(\mathcal E\). We regard them as one-object internal categories
\[
    \mathbb M=(1,M,!,!,e_M,m_{\mathbb M}),
    \qquad
    \mathbb N=(1,N,!,!,e_N,m_{\mathbb N}),
\]
where
\[
    m_{\mathbb M}(a,b)=ba,
    \qquad
    m_{\mathbb N}(\alpha,\beta)=\beta\alpha.
\]
Equivalently, define
\[
    \rho
    :=
    \delta\circ\tau_{P,M}
    :
    P\times M\to P,
\]
where
\[
    \tau_{P,M}:P\times M\to M\times P
\]
is the symmetry. We write
\[
    x\cdot m
    :=
    \rho(x,m)
    =
    \delta(m,x).
\]
The state-transition unit and multiplication laws are then precisely the
internal right-action laws
\[
    x\cdot 1=x,
    \qquad
    x\cdot(mn)=(x\cdot m)\cdot n.
\]
Equivalently, as morphism equations,
\[
    \rho\circ\langle 1_P,e_M!\rangle
    =
    1_P
\]
and
\[
    \rho\circ(\rho\times 1_M)
    =
    \rho\circ(1_P\times\cdot_M)
\]
as maps
\[
    P\times M\times M\to P.
\]

The output map \(\omega\) is the corresponding internal \(N\)-valued output
cocycle:
\[
    \omega(1,x)=1,
    \qquad
    \omega(mn,x)=\omega(m,x)\omega(n,x\cdot m).
\]
The sequential execution picture is
\[
    \begin{tikzcd}[column sep=large]
            x
                \arrow[r, "{m/\omega(m,x)}"]
        &
            x\cdot m
                \arrow[r, "{n/\omega(n,x\cdot m)}"]
        &
            (x\cdot m)\cdot n.
    \end{tikzcd}
\]

\begin{proposition}[One-object Mealy morphisms]
\label{prop:one-object-mealy-morphisms}
    Let \(M\) and \(N\) be monoid objects in \(\mathcal E\), regarded as one-object internal categories by the convention
    \[
        b\circ a=ba.
    \]
    To give a one-object Mealy morphism
    \[
        M\rightsquigarrow N
    \]
    is equivalently to give an internal right \(M\)-object
    \[
        \rho:P\times M\to P
    \]
    together with a map
    \[
        \omega:M\times P\to N
    \]
    satisfying
    \[
        \omega(1,x)=1,
        \qquad
        \omega(mn,x)
        =
        \omega(m,x)\omega(n,x\cdot m),
    \]
    where
    \[
        x\cdot m:=\rho(x,m).
    \]
\end{proposition}
\begin{proof}
    \leavevmode
    \begin{enumerate}
        \item \textbf{Identify the carrier and transition map.} In the one-object case, the carrier span is simply an object \(P\), and the pullbacks defining \(P\circ M\) and \(N\circ P\) reduce to products. A Mealy structure therefore consists of maps
        \[
            \delta:M\times P\to P,
            \qquad
            \omega:M\times P\to N.
        \]
        Using the cartesian symmetry, define
        \[
            \rho
            :=
            \delta\circ\tau_{P,M}
            :
            P\times M\to P.
        \]
        Conversely,
        \[
            \delta
            =
            \rho\circ\tau_{M,P}.
        \]
        Hence specifying \(\delta\) is equivalent to specifying the conventionally typed operation \(\rho\).

        \item \textbf{Identify the right-action laws.} Proposition~\ref{prop:mealy-laws}, together with the convention \(b\circ a=ba\), gives
        \[
            \delta(1,x)=x
        \]
        and
        \[
            \delta(mn,x)
            =
            \delta(n,\delta(m,x)).
        \]
        In terms of
        \[
            x\cdot m
            :=
            \rho(x,m)
            =
            \delta(m,x),
        \]
        these equations become
        \[
            x\cdot1=x
        \]
        and
        \[
            x\cdot(mn)
            =
            (x\cdot m)\cdot n.
        \]
        Thus \(\rho\) is an internal right action of \(M\) on \(P\).

        \item \textbf{Identify the output-cocycle laws.} The Mealy unit law for the output is
        \[
            \omega(1,x)=1.
        \]
        Its multiplication law is
        \[
            \omega(mn,x)
            =
            \omega(m,x)\omega(n,\delta(m,x)),
        \]
        which, using
        \[
            \delta(m,x)=x\cdot m,
        \]
        is exactly
        \[
            \omega(mn,x)
            =
            \omega(m,x)\omega(n,x\cdot m).
        \]

        \item \textbf{Reverse the construction.} Conversely, let
        \[
            \rho:P\times M\to P
        \]
        be an internal right action, and let
        \[
            \omega:M\times P\to N
        \]
        satisfy the displayed unit and cocycle equations. Define
        \[
            \delta
            :=
            \rho\circ\tau_{M,P}
            :
            M\times P\to P.
        \]
        The right-action equations give the transition unit and multiplication laws, while the equations for \(\omega\) give the output unit and multiplication laws. Hence \((P,\delta,\omega)\) determines a one-object Mealy morphism.
    \end{enumerate}
\end{proof}

\begin{remark}[Set-based one-object systems]
\label{rem:set-one-object-systems}
    When \(\mathcal E=\mathbf{Set}\), Proposition~\ref{prop:one-object-mealy-morphisms} says that a one-object Mealy morphism is a set with a right action of the input monoid, together with an output cocycle valued in the output monoid. The set-theoretic reading is therefore a specialization of the internal statement, not part of the definition.
\end{remark}

\subsection{The stateless specialization}
\label{subsec:stateless-mealy-morphisms}

A Mealy morphism
\[
    (P,\Phi):\mathbb A\rightsquigarrow\mathbb B
\]
has a state object \(P\). The map-representable case is obtained by taking the canonical carrier
\[
    F_{0!}
    =
    \left(
        A_0
        \xleftarrow{1_{A_0}}
        A_0
        \xrightarrow{F_0}
        B_0
    \right).
\]
The state is then the current input object:
\[
    \begin{tikzcd}
        & 
            A_0 
                \arrow[dl, "1_{A_0}"'] 
                \arrow[dr, "F_0"] 
        & 
        \\
            A_0 
        &
        & 
            B_0.
    \end{tikzcd}
\]

Let
\[
    F_0:A_0\to B_0
\]
be a morphism in \(\mathcal E\). A Mealy structure on \(F_{0!}\) is a 2-cell
\[
    \Phi:F_{0!}\circ A\Rightarrow B\circ F_{0!}.
\]
First,
\[
    F_{0!}\circ A
    \cong
    \left(
        A_0
        \xleftarrow{s_A}
        A_1
        \xrightarrow{F_0t_A}
        B_0
    \right).
\]
This comes from the pullback
\[
    \begin{tikzcd}
            A_1\times_{t_A,A_0,1_{A_0}}A_0
                \arrow[r]
                \arrow[d]
                \arrow[dr, phantom, "\lrcorner", very near start]
        &
            A_0 
                \arrow[d, "1_{A_0}"]
        \\
            A_1 
                \arrow[r, "t_A"']
        &
            A_0,
    \end{tikzcd}
\]
which is canonically isomorphic to \(A_1\).

Second,
\[
    B\circ F_{0!}
    =
    \left(
        A_0
        \xleftarrow{\pi_1}
        A_0\times_{F_0,B_0,s_B}B_1
        \xrightarrow{t_B\pi_2}
        B_0
    \right),
\]
where
\[
    \begin{tikzcd}
            A_0\times_{F_0,B_0,s_B}B_1
                \arrow[r, "\pi_2"]
                \arrow[d, "\pi_1"']
                \arrow[dr, phantom, "\lrcorner", very near start]
        &
            B_1 
                \arrow[d, "s_B"]
        \\
            A_0 
                \arrow[r, "F_0"']
        &
            B_0.
    \end{tikzcd}
\]
Thus \(\Phi\) is equivalently a map
\[
    \Phi:A_1\to A_0\times_{B_0}B_1.
\]
Write
\[
    \Phi(a)=\bigl(\delta(a),F_1(a)\bigr).
\]
The condition that \(\Phi(a)\) lies in the pullback says
\[
    F_0(\delta(a))=s_B(F_1(a)).
\]
The condition that \(\Phi\) is a map of spans says
\[
    \delta(a)=s_A(a),
    \qquad
    t_B(F_1(a))=F_0(t_A(a)).
\]
Therefore
\[
    s_BF_1=F_0s_A,
    \qquad
    t_BF_1=F_0t_A.
\]
The arrow part is summarized by
\[
    \begin{tikzcd}[column sep=large]
            A_1
                \arrow[r, "F_1"]
                \arrow[d, "{(s_A,t_A)}"']
        &
            B_1
                \arrow[d, "{(s_B,t_B)}"]
        \\
            A_0\times A_0
                \arrow[r, "F_0\times F_0"']
        &
            B_0\times B_0.
    \end{tikzcd}
\]

\begin{definition}[Internal functor]
\label{def:internal-functor}
    Let
    \[
        \mathbb A=(A_0,A_1,s_A,t_A,e_A,m_A)
    \]
    and
    \[
        \mathbb B=(B_0,B_1,s_B,t_B,e_B,m_B)
    \]
    be internal categories in \(\mathcal E\). An \textbf{internal functor}
    \[
        F:\mathbb A\to\mathbb B
    \]
    consists of maps
    \[
        F_0:A_0\to B_0,
        \qquad
        F_1:A_1\to B_1,
    \]
    such that
    \begin{align}
        s_BF_1&=F_0s_A,
        \label{eq:internal-functor-source}
        \\
        t_BF_1&=F_0t_A,
        \label{eq:internal-functor-target}
        \\
        F_1e_A&=e_BF_0,
        \label{eq:internal-functor-identity}
        \\
        F_1m_A&=m_BF_1^{(2)}.
        \label{eq:internal-functor-composition}
    \end{align}
    Here
    \[
        F_1^{(2)}:
        A_1\times_{t_A,A_0,s_A}A_1
        \longrightarrow
        B_1\times_{t_B,B_0,s_B}B_1
    \]
    is the unique map induced by the pullback defining the object of composable pairs in \(\mathbb B\), with projections
    \[
        \pi_1F_1^{(2)}=F_1\pi_1,
        \qquad
        \pi_2F_1^{(2)}=F_1\pi_2.
    \]
    It is well-defined because, on
    \[
        A_1\times_{t_A,A_0,s_A}A_1,
    \]
    one has
    \[
        t_BF_1\pi_1
        =
        F_0t_A\pi_1
        =
        F_0s_A\pi_2
        =
        s_BF_1\pi_2.
    \]

    The source-target equations say that
    \[
        \begin{tikzcd}[column sep=large]
                A_1
                    \arrow[r, "F_1"]
                    \arrow[d, "s_A"']
            &
                B_1
                    \arrow[d, "s_B"]
            \\
                A_0
                    \arrow[r, "F_0"']
            &
                B_0
        \end{tikzcd}
        \qquad
        \begin{tikzcd}[column sep=large]
                A_1
                    \arrow[r, "F_1"]
                    \arrow[d, "t_A"']
            &
                B_1
                    \arrow[d, "t_B"]
            \\
                A_0
                    \arrow[r, "F_0"']
            &
                B_0
        \end{tikzcd}
    \]
    commute, while the identity and composition equations say that
    \[
        \begin{tikzcd}[column sep=large]
                A_0
                    \arrow[r, "F_0"]
                    \arrow[d, "e_A"']
            &
                B_0
                    \arrow[d, "e_B"]
            \\
                A_1
                    \arrow[r, "F_1"']
            &
                B_1
        \end{tikzcd}
    \]
    commutes and that \(F_1\) preserves composition in the sense that the diagram
    \[
        \begin{tikzcd}[column sep=large]
                A_1\times_{t_A,A_0,s_A}A_1
                    \arrow[r, "F_1^{(2)}"]
                    \arrow[d, "m_A"']
            &
                B_1\times_{t_B,B_0,s_B}B_1
                    \arrow[d, "m_B"]
            \\
                A_1
                    \arrow[r, "F_1"']
            &
                B_1
        \end{tikzcd}
    \]
    commutes.
\end{definition}

\begin{theorem}[Stateless Mealy morphisms are internal functors]
\label{thm:representable-mealy-internal-functor}
    Let
    \[
        \mathbb A,\mathbb B
    \]
    be internal categories in \(\mathcal E\), and let
    \[
        F_0:A_0\to B_0
    \]
    be a morphism in \(\mathcal E\). To give an internal functor
    \[
        F:\mathbb A\to\mathbb B
    \]
    with object map \(F_0\) is equivalently to give a Mealy morphism
    \[
        (F_{0!},\Phi):\mathbb A\rightsquigarrow\mathbb B
    \]
    whose carrier is the canonical map-representable span
    \[
        F_{0!}
        =
        \left(
            A_0
            \xleftarrow{1_{A_0}}
            A_0
            \xrightarrow{F_0}
            B_0
        \right).
    \]
\end{theorem}
\begin{proof}
    \leavevmode
    \begin{enumerate}
        \item \textbf{Identify the structure cell.} By the calculation preceding the theorem, a 2-cell
        \[
            \Phi:F_{0!}\circ A\Rightarrow B\circ F_{0!}
        \]
        is equivalently a map
        \[
            \Phi:A_1\to A_0\times_{F_0,B_0,s_B}B_1.
        \]
        Writing \(\Phi(a)=(s_A(a),F_1(a))\), the span-map and pullback conditions are exactly
        \[
            s_BF_1=F_0s_A,
            \qquad
            t_BF_1=F_0t_A.
        \]

        \item \textbf{Identify the laws.} Under this identification, Proposition~\ref{prop:mealy-laws} says that the Mealy unit law is
        \[
            F_1e_A=e_BF_0,
        \]
        and the Mealy multiplication law is
        \[
            F_1(b\circ a)=F_1(b)\circ F_1(a),
        \]
        equivalently
        \[
            F_1m_A=m_BF_1^{(2)}.
        \]
        Hence a stateless Mealy structure on \(F_{0!}\) is exactly an internal functor with object map \(F_0\).

        \item \textbf{Construct the inverse assignment.} Conversely, an internal functor \((F_0,F_1)\) defines
        \[
            \Phi(a)=(s_A(a),F_1(a)).
        \]
        The source-target equations for \(F_1\) make this a map of spans, and identity and composition preservation are exactly the Mealy unit and multiplication laws. The two constructions are inverse.
    \end{enumerate}
\end{proof}

\begin{corollary}[Map-representable Mealy morphisms]
\label{cor:map-representable-mealy-internal-functor-up-to-iso}
    Let
    \[
        (P,\Phi):\mathbb A\rightsquigarrow\mathbb B
    \]
    be a Mealy morphism whose carrier span
    \[
        A_0\xleftarrow{p}P\xrightarrow{q}B_0
    \]
    is map-representable. Put
    \[
        F_0:=q p^{-1}:A_0\to B_0.
    \]
    Then the span isomorphism
    \[
        \sigma:=p:
        \left(
            A_0\xleftarrow{p}P\xrightarrow{q}B_0
        \right)
        \overset{\cong}{\Longrightarrow}
        \left(
            A_0\xleftarrow{1_{A_0}}A_0\xrightarrow{F_0}B_0
        \right)
        =
        F_{0!}
    \]
    transports \(\Phi\) to a Mealy structure
    \[
        \Phi^F:F_{0!}\circ A\Rightarrow B\circ F_{0!}.
    \]
    Consequently, \((P,\Phi)\) determines an internal functor
    \[
        F:\mathbb A\to\mathbb B
    \]
    with object map \(F_0\).
\end{corollary}
\begin{proof}
    \leavevmode
    \begin{enumerate}
        \item \textbf{Construct the carrier-span isomorphism.} Since the carrier is map-representable, the left leg
        \[
            p:P\to A_0
        \]
        is an isomorphism. Put
        \[
            F_0=qp^{-1}.
        \]
        The map
        \[
            \sigma:=p:P\to A_0
        \]
        is then a map of spans
        \[
            \sigma:
            \left(
                A_0\xleftarrow{p}P\xrightarrow{q}B_0
            \right)
            \Longrightarrow
            \left(
                A_0\xleftarrow{1_{A_0}}A_0\xrightarrow{F_0}B_0
            \right),
        \]
        because
        \[
            1_{A_0}\sigma=p
        \]
        and
        \[
            F_0\sigma
            =
            qp^{-1}p
            =
            q.
        \]
        Its inverse as a span map is \(p^{-1}:A_0\to P\).

        \item \textbf{Transport the Mealy structure.} Define
        \[
            \Phi^F
            :=
            (B\sigma)\circ
            \Phi\circ
            (\sigma^{-1}A).
        \]
        Explicitly, this is the composite
        \[
            F_{0!}\circ A
                \xRightarrow{\;\sigma^{-1}A\;}
            P\circ A
                \xRightarrow{\;\Phi\;}
            B\circ P
                \xRightarrow{\;B\sigma\;}
            B\circ F_{0!}.
        \]
        Since \(\sigma\) is invertible, conjugating the unit and multiplication equations for \(\Phi\) by \(\sigma\) gives the unit and multiplication equations for \(\Phi^F\). Hence
        \[
            (F_{0!},\Phi^F):\mathbb A\rightsquigarrow\mathbb B
        \]
        is a Mealy morphism.

        \item \textbf{Apply the stateless correspondence.} By Theorem~\ref{thm:representable-mealy-internal-functor}, a Mealy structure on the canonical carrier \(F_{0!}\) is equivalently an internal functor
        \[
            F:\mathbb A\to\mathbb B
        \]
        whose object map is \(F_0\). Therefore the original map-representable Mealy morphism determines such an internal functor.
    \end{enumerate}
\end{proof}

\section{Mealy simulations}
\label{sec:mealy-simulations}
\subsection{From strict carrier maps to output-comparison cells}
\label{subsec:strict-comparisons-output-comparisons}

The previous section identified the 1-cells. The next question is how to compare two such oriented structures.

Let
\[
    (P,\Phi),(Q,\Psi):\mathbb A\rightsquigarrow\mathbb B
\]
be Mealy morphisms, with carrier spans
\[
    A_0\xleftarrow{p_P}P\xrightarrow{q_P}B_0,
    \qquad
    A_0\xleftarrow{p_Q}Q\xrightarrow{q_Q}B_0.
\]
The first comparison available in the span bicategory is a map of carrier spans. In the direction used below, this is a map
\[
    \theta:Q\to P
\]
satisfying
\[
    p_P\theta=p_Q,
    \qquad
    q_P\theta=q_Q.
\]
It is displayed as
\[
    \begin{tikzcd}[column sep=large]
        &
            Q
                \arrow[d, "\theta" description]
                \arrow[dl, "p_Q"']
                \arrow[dr, "q_Q"]
        &
        \\
            A_0
        &
            P
                \arrow[l, "p_P"]
                \arrow[r, "q_P"']
        &
            B_0.
    \end{tikzcd}
\]
Requiring this carrier map to preserve the Mealy structures gives the equations
\[
    \theta(\delta_Q(a,y))
    =
    \delta_P(a,\theta y),
    \qquad
    \omega_Q(a,y)
    =
    \omega_P(a,\theta y).
\]
Thus a strict comparison preserves the input interface, the output interface, the state transition, and the emitted output arrow.

This gives a valid strict output-preserving comparison. Its behaviour in the map-representable case, however, is too rigid. Let
\[
    F,G:\mathbb A\to\mathbb B
\]
be internal functors, regarded as Mealy morphisms with canonical carriers
\[
    F_{0!}
    =
    \left(
        A_0
        \xleftarrow{1_{A_0}}
        A_0
        \xrightarrow{F_0}
        B_0
    \right),
    \qquad
    G_{0!}
    =
    \left(
        A_0
        \xleftarrow{1_{A_0}}
        A_0
        \xrightarrow{G_0}
        B_0
    \right).
\]
A strict carrier map
\[
    G_{0!}\to F_{0!}
\]
is a map
\[
    \theta:A_0\to A_0
\]
such that
\[
    1_{A_0}\theta=1_{A_0},
    \qquad
    F_0\theta=G_0.
\]
Hence
\[
    \theta=1_{A_0},
    \qquad
    F_0=G_0.
\]
Thus, already at the level of carrier spans, a strict comparison can exist
only when the two functors have the same object map.

If the carrier map is required to preserve the Mealy structures strictly, then
the output-preservation equation gives, for every arrow \(a\in A_1\),
\[
    G_1(a)
    =
    F_1(a).
\]
Indeed, the only possible carrier map is \(1_{A_0}\), and the strict output
equation specializes to
\[
    \omega_G(a,t_Aa)
    =
    \omega_F(a,t_Aa),
\]
that is,
\[
    G_1(a)=F_1(a).
\]
Consequently a strict comparison between canonical map-representable Mealy
morphisms exists only when
\[
    F=G,
\]
and in that case it is the identity strict comparison.

Strict carrier comparisons therefore do not recover nonidentity internal
natural transformations. The component arrow
\[
    \alpha_x:F_0x\to G_0x
\]
has no place in the strict carrier-map datum, even when the two functors are
related by a nontrivial internal natural transformation.

This calculation indicates the required enlargement. A comparison from \((P,\Phi)\) to \((Q,\Psi)\) should still send each \(Q\)-state to a \(P\)-state over the same input interface:
\[
    \theta:Q\to P,
    \qquad
    p_P\theta=p_Q.
\]
The equality of output interfaces should be replaced by an arrow in the output category
\[
    \alpha_y:q_P(\theta y)\to q_Q(y).
\]
In the stateless case this is exactly the shape of a natural-transformation component.

The span calculus already contains the object that stores such data. The composite
\[
    B\circ P
\]
has apex
\[
    P\times_{q_P,B_0,s_B}B_1
\]
and is the span
\[
    A_0
    \xleftarrow{p_P\pi_P}
    P\times_{q_P,B_0,s_B}B_1
    \xrightarrow{t_B\pi_{B_1}}
    B_0.
\]
The defining pullback is
\[
    \begin{tikzcd}
        P\times_{q_P,B_0,s_B}B_1
            \arrow[r, "\pi_{B_1}"]
            \arrow[d, "\pi_P"']
            \arrow[dr, phantom, "\lrcorner", very near start]
        &
            B_1
                \arrow[d, "s_B"]
        \\
            P
                \arrow[r, "q_P"']
        &
            B_0.
    \end{tikzcd}
\]
A generalized element of \(B\circ P\) is a pair
\[
    (x,\beta)
\]
where
\[
    x\in P,
    \qquad
    \beta:q_P(x)\to b
\]
is an arrow of \(\mathbb B\). Therefore a comparison assigning to a \(Q\)-state a \(P\)-state together with an output-comparison arrow is exactly a map of spans
\[
    \Theta:Q\Rightarrow B\circ P.
\]
It is displayed as
\[
    \begin{tikzcd}[column sep=large]
        &
            Q
                \arrow[d, "\Theta" description]
                \arrow[dl, "p_Q"']
                \arrow[dr, "q_Q"]
        &
        \\
            A_0
        &
            P\times_{B_0}B_1
                \arrow[l, "p_P\pi_P"]
                \arrow[r, "t_B\pi_{B_1}"']
        &
            B_0.
    \end{tikzcd}
\]
Writing
\[
    \Theta(y)=\bigl(\theta(y),\alpha_y\bigr),
\]
the span-map and pullback conditions give
\[
    p_P\theta=p_Q,
    \qquad
    s_B\alpha=q_P\theta,
    \qquad
    t_B\alpha=q_Q.
\]
Thus each state \(y\in Q\) is compared with a state \(\theta y\in P\) over the same input interface, and the output interfaces are compared by an arrow
\[
    \begin{tikzcd}[column sep=large]
            q_P(\theta y)
                \arrow[r, "\alpha_y"]
        &
            q_Q(y)
    \end{tikzcd}
\]
in \(\mathbb B\).

For canonical map-representable carriers, a map of spans
\[
    G_{0!}\Rightarrow B\circ F_{0!}
\]
is equivalently a map
\[
    \alpha:A_0\to B_1
\]
satisfying
\[
    s_B\alpha=F_0,
    \qquad
    t_B\alpha=G_0.
\]
These are precisely the component-typing equations for an internal natural transformation
\[
    \alpha:F\Rightarrow G.
\]
The remaining compatibility condition should therefore be the internal naturality equation
\[
    \alpha_v\circ F_1(a)=G_1(a)\circ\alpha_u
\]
for every arrow \(a:u\to v\) of \(\mathbb A\).

Now notice that on the local span category
\[
    \mathcal C_{A,B}=\mathbf{Span}(\mathcal E)(A_0,B_0),
\]
postcomposition by \(B\) defines an endofunctor
\[
    T_B=B\circ-.
\]
Its unit attaches an identity output-comparison arrow, and its multiplication composes output-comparison arrows. Thus \(T_B\) is the output-comparison monad.
A map
\[
    Q\to T_B(P)
\]
is exactly the comparison cell just described.

In the case of the source category \(\mathbb A\), we also have the supplied the input-processing monad
\[
    R_A=-\circ A.
\]
The associativity isomorphisms of the span bicategory identify
\[
    R_AT_B(P)=(B\circ P)\circ A
\]
with
\[
    T_BR_A(P)=B\circ(P\circ A),
\]
and give a distributive law
\[
    R_AT_B\Rightarrow T_BR_A.
\]
Consequently \(R_A\) lifts to a monad on the Kleisli category of \(T_B\).

With this organization, a Mealy morphism is an algebra for the lifted input-processing monad, and a Mealy simulation is an oppositely displayed morphism of such algebras. The displayed direction is chosen so that, in the map-representable case, a simulation
\[
    F_{0!}\Rightarrow G_{0!}
\]
has component arrows
\[
    F_0x\to G_0x
\]
and hence recovers an internal natural transformation
\[
    F\Rightarrow G.
\]

\subsection{The output-comparison monad}
\label{subsec:output-comparison-monad}

Fix internal categories
\[
    \mathbb A=(A_0,A_1,s_A,t_A,e_A,m_A),
    \qquad
    \mathbb B=(B_0,B_1,s_B,t_B,e_B,m_B).
\]
Let
\[
    \mathcal C_{A,B}:=\mathbf{Span}(\mathcal E)(A_0,B_0)
\]
be the local span category. Its objects are spans
\[
    A_0\xleftarrow{p}P\xrightarrow{q}B_0,
\]
and its morphisms are maps of spans. Thus a morphism
\[
    f:P\to Q
\]
in \(\mathcal C_{A,B}\) is a commuting diagram
\[
    \begin{tikzcd}[column sep=large]
        & 
            P 
                \arrow[d, "f" description] 
                \arrow[dl, "p_P"'] 
                \arrow[dr, "q_P"] 
        & 
        \\
            A_0 
        & 
            Q 
                \arrow[l, "p_Q"] 
                \arrow[r, "q_Q"'] 
        & 
            B_0.
    \end{tikzcd}
\]
Equivalently,
\[
    p_Qf=p_P,
    \qquad
    q_Qf=q_P.
\]

\begin{lemma}[Local monads induced by endpoint monads]
\label{lem:local-monads-induced-by-endpoint-monads}
    Let \(A\) be a monad on an object \(A_0\), and let \(B\) be a monad on an
    object \(B_0\), in a bicategory \(\mathcal B\). Write
    \[
        a_{H,G,F}:(H\circ G)\circ F
        \Longrightarrow
        H\circ(G\circ F)
    \]
    for the associator, and write
    \[
        \ell_F:1\circ F\Longrightarrow F,
        \qquad
        r_F:F\circ1\Longrightarrow F
    \]
    for the unitors.

    On the hom-category
    \[
        \mathcal B(A_0,B_0),
    \]
    precomposition by \(A\) defines a monad
    \[
        R=-\circ A,
    \]
    and postcomposition by \(B\) defines a monad
    \[
        T=B\circ-.
    \]

    For a 1-cell \(P:A_0\to B_0\), the unit and multiplication of \(R\) are
    \[
        \eta^R_P
        :=
        (P\eta_A)\circ r_P^{-1}
        :
        P
        \Longrightarrow
        P\circ A,
    \]
    and
    \[
        \mu^R_P
        :=
        (P\mu_A)\circ a_{P,A,A}
        :
        (P\circ A)\circ A
        \Longrightarrow
        P\circ A.
    \]

    The unit and multiplication of \(T\) are
    \[
        \eta^T_P
        :=
        (\eta_B P)\circ\ell_P^{-1}
        :
        P
        \Longrightarrow
        B\circ P,
    \]
    and
    \[
        \mu^T_P
        :=
        (\mu_B P)\circ a_{B,B,P}^{-1}
        :
        B\circ(B\circ P)
        \Longrightarrow
        B\circ P.
    \]

    On a 2-cell \(f:P\Rightarrow Q\), the two endofunctors act by whiskering:
    \[
        R(f)=fA,
        \qquad
        T(f)=Bf.
    \]
\end{lemma}
\begin{proof}
    Naturality of the displayed units and multiplications follows from
    naturality of the bicategorical associators and unitors together with
    functoriality of whiskering.

    For \(R\), the two unit equations reduce, after expanding
    \(\eta^R\) and \(\mu^R\), to the two unit equations for the monad \(A\),
    with the triangle identity supplying the required cancellation of
    unitors. The multiplication equation reduces to the multiplication
    equation for \(A\), with the pentagon identity identifying the two
    reassociations of
    \[
        ((P\circ A)\circ A)\circ A.
    \]

    Dually, the two unit equations for \(T\) reduce to the unit equations for
    \(B\), and the multiplication equation reduces to the multiplication
    equation for \(B\). The required changes of parentheses are supplied by
    the triangle and pentagon identities.

    Hence \(R=-\circ A\) and \(T=B\circ-\) are monads on
    \(\mathcal B(A_0,B_0)\).
\end{proof}

\begin{definition}[The output-comparison endofunctor]
\label{def:output-comparison-endofunctor}
    Define an endofunctor
    \[
        T_B:\mathcal C_{A,B}\to\mathcal C_{A,B}
    \]
    by
    \[
        T_B(P):=B\circ P.
    \]
    Explicitly, if
    \[
        P=
        \left(
            A_0
            \xleftarrow{p}
            P
            \xrightarrow{q}
            B_0
        \right),
    \]
    then \(T_B(P)\) is obtained from the pullback
    \[
        \begin{tikzcd}
                P\times_{q,B_0,s_B}B_1
                    \arrow[r, "\pi_{B_1}"]
                    \arrow[d, "\pi_P"']
                    \arrow[dr, phantom, "\lrcorner", very near start]
            &
                B_1
                    \arrow[d, "s_B"]
            \\
                P
                    \arrow[r, "q"']
            &
                B_0,
        \end{tikzcd}
    \]
    and is the span
    \[
        T_B(P)
        =
        \left(
            A_0
            \xleftarrow{p\pi_P}
            P\times_{q,B_0,s_B}B_1
            \xrightarrow{t_B\pi_{B_1}}
            B_0
        \right).
    \]
    Thus
    \[
        \begin{tikzcd}[column sep=large]
                A_0
            &
                P\times_{B_0}B_1
                    \arrow[l, "p\pi_P"']
                    \arrow[r, "t_B\pi_{B_1}"]
            &
                B_0
        \end{tikzcd}
    \]
    is the carrier span \(T_B(P)\).

    Let
    \[
        f:P\to Q
    \]
    be a map of carrier spans, so that
    \[
        p_Qf=p_P,
        \qquad
        q_Qf=q_P.
    \]
    Define
    \[
        T_B(f):
        P\times_{q_P,B_0,s_B}B_1
        \longrightarrow
        Q\times_{q_Q,B_0,s_B}B_1
    \]
    to be the unique map induced by the pair
    \[
        f\pi_P:
        P\times_{q_P,B_0,s_B}B_1\to Q,
        \qquad
        \pi_{B_1}:
        P\times_{q_P,B_0,s_B}B_1\to B_1.
    \]
    It is well-defined because
    \[
        q_Qf\pi_P
        =
        q_P\pi_P
        =
        s_B\pi_{B_1}.
    \]
    Equivalently, it is characterized by
    \[
        \pi_QT_B(f)=f\pi_P,
        \qquad
        \pi_{B_1}T_B(f)=\pi_{B_1}.
    \]
    The defining diagram is
    \[
        \begin{tikzcd}[column sep=large, row sep=large]
                P\times_{q_P,B_0,s_B}B_1
                    \arrow[r, "{T_B(f)}"]
                    \arrow[d, "\pi_P"']
            &
                Q\times_{q_Q,B_0,s_B}B_1
                    \arrow[r, "\pi_{B_1}"]
                    \arrow[d, "\pi_Q"]
                    \arrow[dr, phantom, "\lrcorner", very near start]
            &
                B_1
                    \arrow[d, "s_B"]
            \\
                P
                    \arrow[r, "f"']
            &
                Q
                    \arrow[r, "q_Q"']
            &
                B_0.
        \end{tikzcd}
    \]
    In generalized-element notation,
    \[
        T_B(f)(x,\beta)=(fx,\beta).
    \]
    The equations characterizing \(T_B(f)\) immediately imply
    \[
        T_B(1_P)=1_{T_B(P)}
    \]
    and
    \[
        T_B(gf)=T_B(g)T_B(f),
    \]
    so \(T_B\) is an endofunctor.
\end{definition}

A generalized element of \(T_B(P)\) is a pair
\[
    (x,\beta)
\]
where \(x\in P\) and
\[
    \beta:q(x)\to b
\]
is an arrow of \(\mathbb B\). Thus \(T_B(P)\) freely attaches an output comparison arrow to a state:
\[
    \begin{tikzcd}[column sep=large]
            q(x) 
                \arrow[r, "\beta"] 
        & 
            b.
    \end{tikzcd}
\]

\begin{proposition}[The output-comparison monad]
\label{prop:output-comparison-monad}
    The endofunctor
    \[
        T_B=B\circ-
    \]
    is a monad on \(\mathcal C_{A,B}\).
\end{proposition}
\begin{proof}
    \leavevmode
    \begin{enumerate}
        \item \textbf{Use the endpoint monad.} This is Lemma~\ref{lem:local-monads-induced-by-endpoint-monads} applied in \(\mathbf{Span}(\mathcal E)\) to the monad \(B\) on \(B_0\). Thus postcomposition by \(B\) gives a monad \(T_B=B\circ-\) on \(\mathcal C_{A,B}\).

        \item \textbf{Unpack the unit.} For a carrier span
        \[
            P=
            \left(
                A_0
                \xleftarrow{p}
                P
                \xrightarrow{q}
                B_0
            \right),
        \]
        the unit is the span map
        \[
            \eta^T_P:P\to T_B(P),
            \qquad
            x\longmapsto (x,e_B(qx)).
        \]
        It is well-defined because \(s_Be_Bq=q\), and its right leg is \(t_Be_Bq=q\).

        \item \textbf{Unpack the multiplication.} A generalized element of \(T_B^2(P)\) is a state \(x\in P\) together with composable arrows
        \[
            \begin{tikzcd}[column sep=large]
                    q(x) 
                        \arrow[r, "\beta"] 
                & 
                    b 
                        \arrow[r, "\gamma"] 
                & 
                    c.
            \end{tikzcd}
        \]
        The multiplication is
        \[
            \mu^T_P(x,\beta,\gamma)=(x,\gamma\circ\beta),
        \]
        or equivalently \((x,m_B(\beta,\gamma))\).

        \item \textbf{Monad laws.} The monad unit and associativity axioms are exactly the identity and associativity laws in the internal category \(\mathbb B\), after whiskering by \(P\). Hence \(T_B\) is a monad.
    \end{enumerate}
\end{proof}

\begin{definition}[The Kleisli category of output comparisons]
\label{def:kleisli-output-category}
    The \textbf{output Kleisli category} associated to \(\mathbb B\) over \((A_0,B_0)\) is the Kleisli category
    \[
        \operatorname{Kl}(T_B).
    \]
    Its objects are carrier spans
    \[
        A_0\xleftarrow{}P\xrightarrow{}B_0.
    \]
    A Kleisli arrow
    \[
        Q\to P
    \]
    is a map in \(\mathcal C_{A,B}\)
    \[
        Q\to T_B(P)=B\circ P.
    \]
\end{definition}

Unpacked, a Kleisli arrow
\[
    Q\to P
\]
is a pair of maps
\[
    \theta:Q\to P,
    \qquad
    \alpha:Q\to B_1,
\]
satisfying
\begin{align}
    p_P\theta&=p_Q,
    \label{eq:kleisli-arrow-input}
    \\
    s_B\alpha&=q_P\theta,
    \label{eq:kleisli-arrow-source}
    \\
    t_B\alpha&=q_Q.
    \label{eq:kleisli-arrow-target}
\end{align}
Equivalently, it is the map of spans
\[
    \begin{tikzcd}[column sep=large]
        & 
            Q
                \arrow[d, "{\langle \theta,\alpha\rangle}" description]
                \arrow[dl, "p_Q"']
                \arrow[dr, "q_Q"]
        &
        \\
            A_0
        &
            P\times_{B_0}B_1
                \arrow[l, "p_P\pi_P"]
                \arrow[r, "t_B\pi_{B_1}"']
        &
            B_0.
    \end{tikzcd}
\]
Thus a state \(y\in Q\) is sent to a state \(\theta y\in P\) over the same input interface, together with an output-comparison arrow
\[
    \begin{tikzcd}[column sep=large]
            q_P(\theta y) 
                \arrow[r, "\alpha_y"] 
        & 
            q_Q(y).
    \end{tikzcd}
\]

Kleisli composition is output-arrow composition. If
\[
    (\theta,\alpha):Q\to P,
    \qquad
    (\kappa,\beta):R\to Q
\]
are Kleisli arrows, then their composite
\[
    R\to P
\]
is represented by
\[
    R
        \xrightarrow{\langle\kappa,\beta\rangle}
    T_B(Q)
        \xrightarrow{T_B\langle\theta,\alpha\rangle}
    T_B^2(P)
        \xrightarrow{\mu^T_P}
    T_B(P).
\]
In components, it has state map
\[
    \theta\kappa:R\to P
\]
and comparison component
\[
    \gamma_z=\beta_z\circ\alpha_{\kappa z}.
\]
The composite comparison arrow is the path
\[
    \begin{tikzcd}[column sep=large]
            q_P(\theta\kappa z)
                \arrow[r, "\alpha_{\kappa z}"]
        &
            q_Q(\kappa z)
                \arrow[r, "\beta_z"]
        &
            q_R(z).
    \end{tikzcd}
\]
The identity Kleisli arrow on \(P\) is
\[
    x\longmapsto (x,e_B(q_Px)).
\]

\subsection{The input-processing monad and Kleisli algebras}
\label{subsec:input-monad-kleisli-algebras}

The source internal category \(\mathbb A\) induces a second local monad.

\begin{definition}[The input-processing endofunctor]
\label{def:input-processing-endofunctor}
    Define an endofunctor
    \[
        R_A:\mathcal C_{A,B}\to\mathcal C_{A,B}
    \]
    by
    \[
        R_A(P):=P\circ A.
    \]
    Explicitly, if
    \[
        P=
        \left(
            A_0
            \xleftarrow{p}
            P
            \xrightarrow{q}
            B_0
        \right),
    \]
    then \(R_A(P)\) is obtained from the pullback
    \[
        \begin{tikzcd}
                A_1\times_{t_A,A_0,p}P
                    \arrow[r, "\pi_P"]
                    \arrow[d, "\pi_{A_1}"']
                    \arrow[dr, phantom, "\lrcorner", very near start]
            &
                P 
                    \arrow[d, "p"]
            \\
                A_1 
                    \arrow[r, "t_A"']
            &
                A_0,
        \end{tikzcd}
    \]
    and is the span
    \[
        R_A(P)
        =
        \left(
            A_0
            \xleftarrow{s_A\pi_{A_1}}
            A_1\times_{t_A,A_0,p}P
            \xrightarrow{q\pi_P}
            B_0
        \right).
    \]
    Thus
    \[
        \begin{tikzcd}[column sep=large]
                A_0
            &
                A_1\times_{A_0}P
                    \arrow[l, "s_A\pi_{A_1}"']
                    \arrow[r, "q\pi_P"]
            &
                B_0
        \end{tikzcd}
    \]
    is the carrier span \(R_A(P)\).

    If
    \[
        f:P\to Q
    \]
    is a map of spans, then
    \[
        R_A(f):R_A(P)\to R_A(Q)
    \]
    is induced by
    \[
        (a,x)\longmapsto(a,fx).
    \]
    This is well-defined because
    \[
        p_Q(fx)=p_P(x)=t_A(a).
    \]
\end{definition}

A generalized element of \(R_A(P)\) is a pair
\[
    (a,x)
\]
where
\[
    a:u\to v
\]
is an arrow of \(\mathbb A\) and \(p(x)=v\). Thus \(R_A(P)\) is the object of states equipped with an admissible input arrow:
\[
    \begin{tikzcd}[column sep=large]
            u 
                \arrow[r, "a"] 
        & 
            v=p(x).
    \end{tikzcd}
\]

\begin{proposition}[The input-processing monad]
\label{prop:input-processing-monad}
    The endofunctor
    \[
        R_A=-\circ A
    \]
    is a monad on \(\mathcal C_{A,B}\).
\end{proposition}
\begin{proof}
    \leavevmode
    \begin{enumerate}
        \item \textbf{Use the endpoint monad.} This is Lemma~\ref{lem:local-monads-induced-by-endpoint-monads} applied in \(\mathbf{Span}(\mathcal E)\) to the monad \(A\) on \(A_0\). Thus precomposition by \(A\) gives a monad \(R_A=-\circ A\) on \(\mathcal C_{A,B}\).

        \item \textbf{Unpack the unit.} For a carrier span
        \[
            P=
            \left(
                A_0
                \xleftarrow{p}
                P
                \xrightarrow{q}
                B_0
            \right),
        \]
        the unit is the span map
        \[
            \eta^R_P:P\to R_A(P),
            \qquad
            x\longmapsto(e_A(px),x).
        \]
        It is well-defined because \(t_Ae_Ap=p\), and its left leg is \(s_Ae_Ap=p\).

        \item \textbf{Unpack the multiplication.}
        A generalized element of \(R_A^2(P)\) is a state \(x\in P\) together with composable input arrows
        \[
            \begin{tikzcd}[column sep=large]
                    u 
                        \arrow[r, "a"] 
                & 
                    v 
                        \arrow[r, "b"] 
                & 
                    w=p(x).
            \end{tikzcd}
        \]
        The multiplication is
        \[
            \mu^R_P(a,b,x)=(b\circ a,x),
        \]
        or equivalently \((m_A(a,b),x)\).

        \item \textbf{Monad laws.} The monad unit and associativity axioms are exactly the identity and associativity laws in the internal category \(\mathbb A\), after whiskering by \(P\). Hence \(R_A\) is a monad.
    \end{enumerate}
\end{proof}

The two local monads are related by span associativity. For every carrier span \(P\), there is a canonical isomorphism
\[
    R_A T_B(P)
    =
    (B\circ P)\circ A
    \cong
    B\circ(P\circ A)
    =
    T_B R_A(P).
\]
On apexes, both sides are canonically isomorphic to the same object of triples
\[
    (a,x,\beta)
\]
with
\[
    t_A(a)=p(x),
    \qquad
    s_B(\beta)=q(x).
\]
The two pullback presentations are displayed as
\[
    \begin{tikzcd}
            A_1\times_{A_0}(P\times_{B_0}B_1)
                \arrow[r]
                \arrow[d]
                \arrow[dr, phantom, "\lrcorner", very near start]
        &
            P\times_{B_0}B_1
                \arrow[d, "p\pi_P"]
        \\
            A_1
                \arrow[r, "t_A"']
        &
            A_0
    \end{tikzcd}
    \qquad
    \begin{tikzcd}
            (A_1\times_{A_0}P)\times_{B_0}B_1
                \arrow[r]
                \arrow[d]
                \arrow[dr, phantom, "\lrcorner", very near start]
        &
            B_1
                \arrow[d, "s_B"]
        \\
            A_1\times_{A_0}P
                \arrow[r, "q\pi_P"']
        &
            B_0.
    \end{tikzcd}
\]
In generalized-element notation, this isomorphism sends
\[
    (a,(x,\beta))
\]
to
\[
    ((a,x),\beta).
\]

\begin{lemma}[Endpoint distributive law]
\label{lem:endpoint-distributive-law}
    Let \(A\) be a monad on \(A_0\), and let \(B\) be a monad on \(B_0\), in a
    bicategory \(\mathcal B\). On the hom-category
    \[
        \mathcal B(A_0,B_0),
    \]
    write
    \[
        R=-\circ A,
        \qquad
        T=B\circ-
    \]
    for the local monads of
    Lemma~\ref{lem:local-monads-induced-by-endpoint-monads}.

    There is a natural isomorphism
    \[
        \lambda:RT\Rightarrow TR
    \]
    whose component at a 1-cell \(P:A_0\to B_0\) is the bicategorical
    associator
    \[
        \lambda_P
        :=
        a_{B,P,A}:
        (B\circ P)\circ A
        \xRightarrow{\cong}
        B\circ(P\circ A).
    \]

    This natural isomorphism is a distributive law of \(R\) over \(T\). Thus
    its components satisfy
    \begin{align}
        \lambda\circ\eta^R T
        &=
        T\eta^R,
        \label{eq:endpoint-distributive-R-unit}
        \\
        \lambda\circ R\eta^T
        &=
        \eta^T R,
        \label{eq:endpoint-distributive-T-unit}
        \\
        \lambda\circ\mu^R T
        &=
        T\mu^R\circ\lambda R\circ R\lambda,
        \label{eq:endpoint-distributive-R-multiplication}
        \\
        \lambda\circ R\mu^T
        &=
        \mu^T R\circ T\lambda\circ\lambda T.
        \label{eq:endpoint-distributive-T-multiplication}
    \end{align}
    We use the convention that a distributive law
    \[
        RT\Rightarrow TR
    \]
    lifts \(R\) to the Kleisli category of \(T\).
\end{lemma}
\begin{proof}
    Naturality of \(\lambda\) is naturality of the associator
    \(a_{B,P,A}\) in the middle variable \(P\).

    For the \(R\)-unit equation, both sides are canonical composites
    \[
        B\circ P
        \Longrightarrow
        B\circ(P\circ A)
    \]
    obtained by inserting the unit of \(A\). After expanding
    \(\eta^R\), their equality is the triangle identity together with
    naturality of the associator.

    For the \(T\)-unit equation, both sides are canonical composites
    \[
        P\circ A
        \Longrightarrow
        B\circ(P\circ A)
    \]
    obtained by inserting the unit of \(B\). Their equality is again the
    triangle identity together with naturality of the associator.

    For the \(R\)-multiplication equation, both sides are composites
    \[
        ((B\circ P)\circ A)\circ A
        \Longrightarrow
        B\circ(P\circ A)
    \]
    which reassociate the four factors and then apply the multiplication of
    \(A\). Naturality of the associator moves the whiskered multiplication
    through the reassociation, and the pentagon identity identifies the two
    possible reassociation composites. This gives
    \[
        \lambda\circ\mu^R T
        =
        T\mu^R\circ\lambda R\circ R\lambda.
    \]

    Similarly, the \(T\)-multiplication equation compares two composites
    \[
        (B\circ(B\circ P))\circ A
        \Longrightarrow
        B\circ(P\circ A)
    \]
    which reassociate the four factors and then apply the multiplication of
    \(B\). Naturality of the associator and the pentagon identity give
    \[
        \lambda\circ R\mu^T
        =
        \mu^T R\circ T\lambda\circ\lambda T.
    \]

    Hence \(\lambda:RT\Rightarrow TR\) is a distributive law.
\end{proof}

\begin{proposition}[The canonical distributive law]
\label{prop:canonical-distributive-law}
    The associativity constraints of \(\mathbf{Span}(\mathcal E)\) determine a distributive-law isomorphism
    \[
        \lambda:R_A T_B\Rightarrow T_B R_A
    \]
    of the input-processing monad over the output-comparison monad.
\end{proposition}
\begin{proof}
    \leavevmode
    \begin{enumerate}
        \item \textbf{Define the component.} By Lemma~\ref{lem:endpoint-distributive-law}, the component at a carrier span \(P\) is the span associator
        \[
            \lambda_P:
            (B\circ P)\circ A
            \xrightarrow{\cong}
            B\circ(P\circ A).
        \]
        On apexes, both sides are canonically pullback presentations of the object of triples \((a,x,\beta)\) satisfying
        \[
            t_A(a)=p(x),
            \qquad
            s_B(\beta)=q(x),
        \]
        and \(\lambda_P\) sends
        \[
            (a,(x,\beta))
            \longmapsto
            ((a,x),\beta).
        \]

        \item \textbf{Check naturality.} If \(f:P\to Q\) is a map of carrier spans, both composites in the naturality square send \((a,(x,\beta))\) to \(((a,fx),\beta)\). Hence \(\lambda\) is natural.

        \item \textbf{Check the distributive-law axioms.} The compatibility with the two units inserts either an identity input arrow or an identity output arrow, and the associator leaves the other datum unchanged. The compatibility with the two multiplications says that reassociation commutes with composition in \(\mathbb A\) and \(\mathbb B\). Componentwise, for example, the \(T_B\)-multiplication compatibility sends
        \[
            (a,x,\beta,\gamma)
        \]
        on both sides to
        \[
            ((a,x),\gamma\circ\beta),
        \]
        while the \(R_A\)-multiplication compatibility sends
        \[
            (a,b,x,\beta)
        \]
        on both sides to
        \[
            ((b\circ a,x),\beta).
        \]
        These are exactly the distributive-law equations for \(R_AT_B\Rightarrow T_BR_A\).
    \end{enumerate}
\end{proof}

\begin{corollary}[The lifted input-processing monad]
\label{cor:lifted-input-monad}
    The distributive law of Proposition~\ref{prop:canonical-distributive-law} lifts \(R_A\) to a monad
    \[
        \overline R_A
    \]
    on the Kleisli category
    \[
        \operatorname{Kl}(T_B).
    \]
\end{corollary}
\begin{proof}
    \leavevmode
    \begin{enumerate}
        \item \textbf{Apply the Kleisli lifting theorem.} By the standard Kleisli lifting theorem for distributive laws of monads, a distributive law \(RT\Rightarrow TR\) lifts \(R\) to a monad on \(\operatorname{Kl}(T)\). Applying this to \(R_A=-\circ A\), \(T_B=B\circ-\), and Proposition~\ref{prop:canonical-distributive-law} gives the lifted monad \(\overline R_A\).

        \item \textbf{Record the formula.} On objects,
        \[
            \overline R_A(P)=R_A(P).
        \]
        If a Kleisli arrow \(f:P\to Q\) is represented in \(\mathcal C_{A,B}\) by a map \(f:P\to T_B(Q)\), then \(\overline R_A(f)\) is represented by
        \[
            R_A(P)
            \xrightarrow{R_A f}
            R_AT_B(Q)
            \xrightarrow{\lambda_Q}
            T_BR_A(Q).
        \]
        Thus, if \(f(x)=(\theta x,\alpha_x)\), then
        \[
            (a,x)
            \longmapsto
            ((a,\theta x),\alpha_x).
        \]
    \end{enumerate}
\end{proof}

\begin{theorem}[Mealy morphisms as Kleisli algebras]
\label{thm:mealy-morphisms-kleisli-algebras}
    Let \(\mathbb A\) and \(\mathbb B\) be internal categories in
    \(\mathcal E\). To give a Mealy morphism
    \[
        \mathbb A\rightsquigarrow\mathbb B
    \]
    is equivalently to give an algebra for the lifted monad
    \[
        \overline R_A
    \]
    on the Kleisli category
    \[
        \operatorname{Kl}(T_B).
    \]
\end{theorem}
\begin{proof}
    \leavevmode
    \begin{enumerate}
        \item \textbf{Unpack the algebra structure.}
        An algebra for \(\overline R_A\) on an object \(P\) of
        \(\operatorname{Kl}(T_B)\) is a Kleisli arrow
        \[
            \Phi:
            \overline R_A(P)\to P.
        \]
        Since
        \[
            \overline R_A(P)=R_A(P),
        \]
        such a Kleisli arrow is represented in
        \(\mathcal C_{A,B}\) by a map
        \[
            \Phi:
            R_A(P)\to T_B(P).
        \]
        Using
        \[
            R_A(P)=P\circ A,
            \qquad
            T_B(P)=B\circ P,
        \]
        this is exactly a 2-cell
        \[
            \Phi:
            P\circ A\Rightarrow B\circ P
        \]
        in \(\mathbf{Span}(\mathcal E)\).

        \item \textbf{Identify the component data.}
        On apexes, \(\Phi\) is a map
        \[
            A_1\times_{t_A,A_0,p}P
            \longrightarrow
            P\times_{q,B_0,s_B}B_1.
        \]
        Write
        \[
            \Phi(a,x)
            =
            \bigl(\delta(a,x),\omega(a,x)\bigr).
        \]
        The condition that this is a map of spans is exactly
        \[
            p\delta=s_A\pi_{A_1},
            \qquad
            s_B\omega=q\delta,
            \qquad
            t_B\omega=q\pi_P.
        \]
        Thus the algebra structure map is precisely the transition-output datum
        of a Mealy morphism.

        \item \textbf{Compute the algebra unit law.}
        The unit of the lifted monad \(\overline R_A\) is the image in
        \(\operatorname{Kl}(T_B)\) of the unit
        \[
            \eta^R_P:P\to R_A(P).
        \]
        As a map in \(\mathcal C_{A,B}\), this Kleisli arrow is represented by
        \[
            P
            \xrightarrow{\eta^R_P}
            R_A(P)
            \xrightarrow{\eta^T_{R_A(P)}}
            T_B R_A(P).
        \]

        The algebra unit equation in \(\operatorname{Kl}(T_B)\) says that the
        Kleisli composite of this arrow with \(\Phi\) is the identity Kleisli
        arrow on \(P\). Expanding Kleisli composition and using the unit laws
        of \(T_B\), this equation reduces to
        \[
            \Phi\circ\eta^R_P
            =
            \eta^T_P
        \]
        as maps
        \[
            P\to T_B(P).
        \]

        Under the endpoint descriptions of the two monads,
        \[
            \eta^R_P=P\eta_A,
            \qquad
            \eta^T_P=\eta_B P.
        \]
        Hence the algebra unit law is exactly
        \[
            \Phi\circ(P\eta_A)
            =
            \eta_B P,
        \]
        which is the abstract Mealy unit law
        \eqref{eq:mealy-unit-abstract}.

        \item \textbf{Compute the algebra multiplication law.}
        The multiplication of the lifted monad is the image in
        \(\operatorname{Kl}(T_B)\) of
        \[
            \mu^R_P:
            R_A^2(P)\to R_A(P).
        \]
        The algebra multiplication equation in the Kleisli category compares
        the two Kleisli arrows
        \[
            R_A^2(P)\to P
        \]
        obtained by first using \(\mu^R_P\), or by first applying the lifted
        endofunctor to \(\Phi\).

        Expanding both Kleisli composites gives the equality
        \[
            \Phi\circ\mu^R_P
            =
            \mu^T_P\circ
            T_B(\Phi)\circ
            \lambda_P\circ
            R_A(\Phi)
        \]
        as maps
        \[
            R_A^2(P)\to T_B(P).
        \]

        Here
        \[
            R_A(\Phi):
            R_A^2(P)\to R_AT_B(P)
        \]
        is followed by the distributive-law component
        \[
            \lambda_P:
            R_AT_B(P)\to T_BR_A(P),
        \]
        then by
        \[
            T_B(\Phi):
            T_BR_A(P)\to T_B^2(P),
        \]
        and finally by
        \[
            \mu^T_P:
            T_B^2(P)\to T_B(P).
        \]

        Under the endpoint descriptions,
        \[
            \mu^R_P=P\mu_A,
            \qquad
            \mu^T_P=\mu_B P,
        \]
        \[
            R_A(\Phi)=\Phi A,
            \qquad
            T_B(\Phi)=B\Phi,
        \]
        while \(\lambda_P\) is the canonical reassociation
        \[
            (B\circ P)\circ A
            \xRightarrow{\cong}
            B\circ(P\circ A).
        \]
        Therefore the algebra multiplication law is exactly
        \[
            \Phi\circ(P\mu_A)
            =
            (\mu_B P)\circ
            (B\Phi)\circ
            \lambda_P\circ
            (\Phi A).
        \]
        With the canonical associativity constraints suppressed, this is
        \[
            \Phi\circ(P\mu_A)
            =
            (\mu_B P)\circ(B\Phi)\circ(\Phi A),
        \]
        namely the abstract Mealy multiplication law
        \eqref{eq:mealy-multiplication-abstract}.

        \item \textbf{Conclude.}
        An \(\overline R_A\)-algebra therefore consists exactly of a carrier
        span \(P\), a structure cell
        \[
            \Phi:P\circ A\Rightarrow B\circ P,
        \]
        and the Mealy unit and multiplication equations. These are precisely
        the data and laws of a Mealy morphism
        \[
            \mathbb A\rightsquigarrow\mathbb B.
        \]
        The two constructions are inverse.
    \end{enumerate}
\end{proof}

Let
\[
    (P,\Phi),(Q,\Psi):\mathbb A\rightsquigarrow\mathbb B
\]
be Mealy morphisms. By Theorem~\ref{thm:mealy-morphisms-kleisli-algebras}, they are algebras for \(\overline R_A\) in \(\operatorname{Kl}(T_B)\). The simulation direction used in this chapter is the direction that recovers ordinary internal natural transformations in the stateless case.

The 2-cells used here are not the strict 2-cells between oriented monad morphisms in \(\mathbf{Span}(\mathcal E)\). They are Kleisli output-comparison cells for the monad \(B\circ-\), with the direction reversed so that the stateless case gives ordinary internal natural transformations.

\begin{definition}[Mealy simulation]
\label{def:mealy-simulation}
    Let
    \[
        (P,\Phi),(Q,\Psi):\mathbb A\rightsquigarrow\mathbb B
    \]
    be Mealy morphisms. A \textbf{Mealy simulation}
    \[
        (P,\Phi)\Rightarrow(Q,\Psi)
    \]
    is a morphism
    \[
        (Q,\Psi)\to(P,\Phi)
    \]
    of \(\overline R_A\)-algebras in the Kleisli category
    \[
        \operatorname{Kl}(T_B).
    \]

    Equivalently, it is a map in the local span category
    \[
        \Theta:Q\to T_B(P),
    \]
    or, in bicategorical notation, a map of spans
    \[
        \Theta:Q\Rightarrow B\circ P,
    \]
    satisfying the algebra-morphism equation
    \begin{equation}
        \mu^T_P\circ
        T_B(\Phi)\circ
        \lambda_P\circ
        R_A(\Theta)
        =
        \mu^T_P\circ
        T_B(\Theta)\circ
        \Psi.
        \label{eq:mealy-simulation-kleisli}
    \end{equation}
    Both sides of this equation are maps
    \[
        R_A(Q)\to T_B(P)
    \]
    in \(\mathcal C_{A,B}\).

    In the span-composition notation of the chapter, equation
    \eqref{eq:mealy-simulation-kleisli} is
    \begin{equation}
        (\mu_B P)\circ
        (B\Phi)\circ
        \lambda_P\circ
        (\Theta A)
        =
        (\mu_B P)\circ
        (B\Theta)\circ
        \Psi.
        \label{eq:mealy-simulation-abstract}
    \end{equation}
    Here
    \[
        \lambda_P:
        (B\circ P)\circ A
        \xRightarrow{\cong}
        B\circ(P\circ A)
    \]
    is the canonical distributive-law component induced by the associativity constraint of \(\mathbf{Span}(\mathcal E)\).

    Writing
    \[
        \Theta(y)=\bigl(\theta(y),\alpha(y)\bigr),
    \]
    the underlying comparison data are maps
    \[
        \theta:Q\to P,
        \qquad
        \alpha:Q\to B_1,
    \]
    satisfying
    \begin{align}
        p_P\theta&=p_Q,
        \label{eq:mealy-simulation-input-interface}
        \\
        s_B\alpha&=q_P\theta,
        \label{eq:mealy-simulation-alpha-source}
        \\
        t_B\alpha&=q_Q.
        \label{eq:mealy-simulation-alpha-target}
    \end{align}
    Thus each state \(y\in Q\) is compared with a state \(\theta y\in P\) over the same input interface, together with an output-comparison arrow
    \[
        \alpha_y:q_P(\theta y)\to q_Q(y)
    \]
    in \(\mathbb B\).
\end{definition}

\begin{proposition}[The simulation laws]
\label{prop:mealy-simulation-laws}
    The abstract compatibility equation \eqref{eq:mealy-simulation-abstract} is equivalent to the following component equations:
    \begin{align}
        \theta(\delta_Q(a,y))
        &=
        \delta_P(a,\theta y),
        \label{eq:mealy-simulation-delta}
        \\
        \alpha_y\circ \omega_P(a,\theta y)
        &=
        \omega_Q(a,y)\circ \alpha_{\delta_Q(a,y)}.
        \label{eq:mealy-simulation-omega}
    \end{align}
\end{proposition}
\begin{proof}
    \leavevmode
    \begin{enumerate}
        \item \textbf{Fix the relevant input data.} The relevant generalized elements are pairs
        \[
            (a,y)
        \]
        satisfying
        \[
            t_A(a)=p_Q(y).
        \]
        Thus \(a:u\to v\) is an input arrow with \(p_Q(y)=v\).

        \item \textbf{Use preservation of the input interface.} Since
        \[
            p_P\theta=p_Q,
        \]
        the pair
        \[
            (a,\theta y)
        \]
        is admissible for \(P\):
        \[
            t_A(a)=p_Q(y)=p_P(\theta y).
        \]
        Hence \(P\) can execute the same input arrow \(a\) at the compared
        state \(\theta y\).

        \item \textbf{Compute the left-hand side of the abstract equation.} The left-hand side of \eqref{eq:mealy-simulation-abstract} first compares \(y\) with \(\theta y\), and then executes \(P\) along \(a\). The state part is
        \[
            \delta_P(a,\theta y).
        \]
        The output arrows form the composite
        \[
            \begin{tikzcd}[column sep=huge]
                    {q_P(\delta_P(a,\theta y))}
                        \arrow[r, "{\omega_P(a,\theta y)}"]
                &
                    {q_P(\theta y)}
                        \arrow[r, "{\alpha_y}"]
                &
                    {q_Q(y)}.
            \end{tikzcd}
        \]
        Therefore the output component is
        \[
            \alpha_y\circ\omega_P(a,\theta y).
        \]

        \item \textbf{Compute the right-hand side of the abstract equation.} The right-hand side first executes \(Q\) along \(a\):
        \[
            y
            \longmapsto
            \delta_Q(a,y),
        \]
        and emits
        \[
            \omega_Q(a,y):q_Q(\delta_Q(a,y))\to q_Q(y).
        \]
        It then compares the updated \(Q\)-state with \(P\), giving
        \[
            \theta(\delta_Q(a,y))
        \]
        and
        \[
            \alpha_{\delta_Q(a,y)}:
            q_P(\theta(\delta_Q(a,y)))\to q_Q(\delta_Q(a,y)).
        \]
        The output arrows form the composite
        \[
            \begin{tikzcd}[column sep=huge]
                    {q_P(\theta(\delta_Q(a,y)))}
                        \arrow[r, "{\alpha_{\delta_Q(a,y)}}"]
                &
                    {q_Q(\delta_Q(a,y))}
                        \arrow[r, "{\omega_Q(a,y)}"]
                &
                    {q_Q(y)}.
            \end{tikzcd}
        \]
        Therefore the output component is
        \[
            \omega_Q(a,y)\circ\alpha_{\delta_Q(a,y)}.
        \]

        \item \textbf{Compare the state components.} Equality of the state components gives
        \[
            \theta(\delta_Q(a,y))
            =
            \delta_P(a,\theta y).
        \]

        \item \textbf{Compare the output components.} Equality of the output components gives
        \[
            \alpha_y\circ \omega_P(a,\theta y)
            =
            \omega_Q(a,y)\circ \alpha_{\delta_Q(a,y)}.
        \]
        Hence the abstract algebra-morphism condition is equivalent to the two displayed simulation laws.
    \end{enumerate}
\end{proof}

\begin{remark}[Strict simulations]
\label{rem:strict-mealy-simulations}
    Suppose
    \[
        \theta:Q\to P
    \]
    satisfies
    \[
        p_P\theta=p_Q,
        \qquad
        q_P\theta=q_Q,
    \]
    and
    \[
        \theta(\delta_Q(a,y))=\delta_P(a,\theta y),
        \qquad
        \omega_Q(a,y)=\omega_P(a,\theta y).
    \]
    Then \(\theta\) defines a Mealy simulation in the sense of Definition~\ref{def:mealy-simulation} by taking
    \[
        \alpha=e_Bq_Q.
    \]
    Conversely, if a Mealy simulation has comparison component
    \[
        \alpha=e_Bq_Q,
    \]
    then the equation
    \[
        s_B\alpha=q_P\theta
    \]
    forces
    \[
        q_P\theta=q_Q,
    \]
    and the simulation laws reduce exactly to strict preservation of state transition and output:
    \[
        \theta(\delta_Q(a,y))=\delta_P(a,\theta y),
        \qquad
        \omega_Q(a,y)=\omega_P(a,\theta y).
    \]
\end{remark}

\subsection{One-object simulations}
\label{subsec:one-object-mealy-simulations}

Let
\[
    M,
    \qquad
    N
\]
be monoid objects in \(\mathcal E\), regarded as one-object internal categories by the convention
\[
    b\circ a=ba
\]
from Subsection~\ref{subsec:two-running-specializations}. Let
\[
    (P,\omega_P),
    \qquad
    (Q,\omega_Q)
\]
be one-object Mealy morphisms
\[
    M\rightsquigarrow N.
\]
Thus \(P\) and \(Q\) are internal right \(M\)-objects, and
\[
    \omega_P:M\times P\to N,
    \qquad
    \omega_Q:M\times Q\to N
\]
are output cocycles.

A Mealy simulation
\[
    P\Rightarrow Q
\]
consists of maps
\[
    \theta:Q\to P,
    \qquad
    \alpha:Q\to N.
\]
The transition equation says that \(\theta\) is \(M\)-equivariant:
\[
    \theta(y\cdot m)=\theta(y)\cdot m.
\]
The output equation says
\[
    \alpha(y)\,\omega_P(m,\theta y)
    =
    \omega_Q(m,y)\,\alpha(y\cdot m).
\]
The corresponding output-comparison diagram is
\[
    \begin{tikzcd}[column sep=huge, row sep=large]
        {}
            \arrow[r, "{\omega_P(m,\theta y)}"]
            \arrow[d, "{\alpha_{y\cdot m}}"']
        &
        {}
            \arrow[d, "{\alpha_y}"]
        \\
        {}
            \arrow[r, "{\omega_Q(m,y)}"']
        &
        {}
    \end{tikzcd}
\]
where the labels multiply according to the one-object convention.Thus one-object Mealy simulations compare output cocycles by a state-dependent cochain. When the output monoid is a group, this becomes the usual nonabelian coboundary relation. The strict output-preserving simulations are the special case
\[
    \alpha(y)=1
\]
for all \(y\).

\begin{proposition}[One-object Mealy simulations]
\label{prop:one-object-mealy-simulations}
    Let
    \[
        P,Q:M\rightsquigarrow N
    \]
    be one-object Mealy morphisms internal to \(\mathcal E\). A Mealy simulation
    \[
        P\Rightarrow Q
    \]
    is equivalently an \(M\)-equivariant map
    \[
        \theta:Q\to P
    \]
    together with a map
    \[
        \alpha:Q\to N
    \]
    satisfying
    \[
        \alpha(y)\,\omega_P(m,\theta y)
        =
        \omega_Q(m,y)\,\alpha(y\cdot m)
    \]
    as a morphism equation \(M\times Q\to N\).
\end{proposition}
\begin{proof}
    \leavevmode
    \begin{enumerate}
        \item \textbf{Specialize the comparison data.} In the one-object case all interface equations are vacuous. Thus a simulation consists of maps
        \[
            \theta:Q\to P,
            \qquad
            \alpha:Q\to N.
        \]

        \item \textbf{Specialize the simulation laws.} Proposition~\ref{prop:mealy-simulation-laws} gives
        \[
            \theta(y\cdot m)=\theta(y)\cdot m,
        \]
        so \(\theta\) is \(M\)-equivariant. Under the one-object convention \(b\circ a=ba\), the output equation becomes
        \[
            \alpha(y)\,\omega_P(m,\theta y)
            =
            \omega_Q(m,y)\,\alpha(y\cdot m).
        \]

        \item \textbf{Reverse the specialization.} Conversely, an \(M\)-equivariant map \(\theta\) and a map \(\alpha\) satisfying the displayed equation determine exactly the component data and laws of a one-object Mealy simulation.
    \end{enumerate}
\end{proof}

The stateless one-object case is obtained by taking \(P=1\). A Mealy morphism
\[
    M\rightsquigarrow N
\]
then reduces to a monoid morphism
\[
    h:M\to N.
\]
A simulation between stateless one-object Mealy morphisms
\[
    h,k:M\to N
\]
is a global element
\[
    \alpha:1\to N
\]
satisfying
\[
    \alpha\,h(m)=k(m)\,\alpha
\]
as a morphism equation \(M\to N\). This is the component form of
\[
    \alpha\circ h(m)=k(m)\circ\alpha
\]
under the one-object convention \(b\circ a=ba\). In \(\mathbf{Set}\), this is precisely a natural transformation between the corresponding one-object functors.

\subsection{Stateless simulations}
\label{subsec:stateless-mealy-simulations}

Let
\[
    F,G:\mathbb A\to\mathbb B
\]
be internal functors, regarded as Mealy morphisms carried by the canonical map-representable spans
\[
    F_{0!}
    =
    \left(
        A_0
        \xleftarrow{1_{A_0}}
        A_0
        \xrightarrow{F_0}
        B_0
    \right),
    \qquad
    G_{0!}
    =
    \left(
        A_0
        \xleftarrow{1_{A_0}}
        A_0
        \xrightarrow{G_0}
        B_0
    \right).
\]
A Mealy simulation
\[
    F_{0!}\Rightarrow G_{0!}
\]
is, by definition, a map of spans
\[
    G_{0!}\Rightarrow B\circ F_{0!}
\]
satisfying the simulation law.

The composite
\[
    B\circ F_{0!}:A_0\to B_0
\]
has apex
\[
    A_0\times_{F_0,B_0,s_B}B_1
\]
and is the span
\[
    A_0
    \xleftarrow{\pi_1}
    A_0\times_{F_0,B_0,s_B}B_1
    \xrightarrow{t_B\pi_2}
    B_0.
\]
The defining pullback is
\[
    \begin{tikzcd}
            A_0\times_{F_0,B_0,s_B}B_1
                \arrow[r, "\pi_2"]
                \arrow[d, "\pi_1"']
                \arrow[dr, phantom, "\lrcorner", very near start]
        &
            B_1 
                \arrow[d, "s_B"]
        \\
            A_0 
                \arrow[r, "F_0"']
        &
            B_0.
    \end{tikzcd}
\]
A map of spans
\[
    \bar\alpha:G_{0!}\Rightarrow B\circ F_{0!}
\]
is therefore equivalently a map
\[
    \bar\alpha:A_0\to A_0\times_{F_0,B_0,s_B}B_1
\]
such that
\[
    \pi_1\bar\alpha=1_{A_0},
    \qquad
    t_B\pi_2\bar\alpha=G_0.
\]
The span-map diagram is
\[
    \begin{tikzcd}[column sep=large]
        & 
            A_0
                \arrow[d, "\bar\alpha" description]
                \arrow[dl, "1_{A_0}"']
                \arrow[dr, "G_0"]
        &
        \\
            A_0
        &
            A_0\times_{B_0}B_1
                \arrow[l, "\pi_1"]
                \arrow[r, "t_B\pi_2"']
        &
            B_0.
    \end{tikzcd}
\]
Writing
\[
    \alpha:=\pi_2\bar\alpha,
\]
the pullback condition gives
\[
    s_B\alpha=F_0,
\]
and the right-leg condition gives
\[
    t_B\alpha=G_0.
\]
Thus a stateless Mealy simulation has components
\[
    \alpha_x:F_0x\to G_0x
\]
in \(\mathbb B\), displayed as
\[
    \begin{tikzcd}[column sep=large]
            F_0x 
                \arrow[r, "\alpha_x"] 
        & 
            G_0x.
    \end{tikzcd}
\]

\begin{definition}[Internal natural transformation]
\label{def:internal-natural-transformation}
    Let
    \[
        F,G:\mathbb A\to\mathbb B
    \]
    be internal functors. An \textbf{internal natural transformation}
    \[
        \alpha:F\Rightarrow G
    \]
    is a map
    \[
        \alpha:A_0\to B_1
    \]
    satisfying
    \[
        s_B\alpha=F_0,
        \qquad
        t_B\alpha=G_0,
    \]
        and
    \[
        m_B\langle F_1,\alpha t_A\rangle
        =
        m_B\langle \alpha s_A,G_1\rangle.
    \]
    Here
    \[
        \langle F_1,\alpha t_A\rangle:
        A_1\to B_1\times_{t_B,B_0,s_B}B_1
    \]
    is well-defined because
    \[
        t_BF_1=F_0t_A=s_B\alpha t_A,
    \]
    and
    \[
        \langle \alpha s_A,G_1\rangle:
        A_1\to B_1\times_{t_B,B_0,s_B}B_1
    \]
    is well-defined because
    \[
        t_B\alpha s_A=G_0s_A=s_BG_1.
    \]
    In generalized-element notation, for an arrow \(a:u\to v\), this equation
    is
    \[
        \alpha_v\circ F_1(a)
        =
        G_1(a)\circ \alpha_u.
    \]
    The naturality square is
    \[
        \begin{tikzcd}[column sep=large, row sep=large]
                F_0u
                    \arrow[r, "F_1(a)"]
                    \arrow[d, "\alpha_u"']
            &
                F_0v
                    \arrow[d, "\alpha_v"]
            \\
                G_0u
                    \arrow[r, "G_1(a)"']
            &
                G_0v.
        \end{tikzcd}
    \]
\end{definition}

\begin{theorem}[Stateless Mealy simulations are internal natural transformations]
\label{thm:stateless-mealy-simulations-natural-transformations}
    Let
    \[
        F,G:\mathbb A\to\mathbb B
    \]
    be internal functors, regarded as canonical map-representable Mealy morphisms. To give a Mealy simulation
    \[
        F_{0!}\Rightarrow G_{0!}
    \]
    is equivalently to give an internal natural transformation
    \[
        F\Rightarrow G.
    \]
\end{theorem}
\begin{proof}
    \leavevmode
    \begin{enumerate}
        \item \textbf{Identify the component map.} A Mealy simulation
        \[
            F_{0!}\Rightarrow G_{0!}
        \]
        is a map of spans
        \[
            G_{0!}\Rightarrow B\circ F_{0!}.
        \]
        By the calculation preceding the theorem, this is equivalently a map
        \[
            \alpha:A_0\to B_1
        \]
        satisfying
        \[
            s_B\alpha=F_0,
            \qquad
            t_B\alpha=G_0.
        \]
        These are exactly the source and target equations for the components of an internal natural transformation.

        \item \textbf{Identify naturality.} For the canonical Mealy structures associated to \(F\) and \(G\), an input arrow \(a:u\to v\) emits \(F_1(a):F_0u\to F_0v\) and \(G_1(a):G_0u\to G_0v\), respectively. Proposition~\ref{prop:mealy-simulation-laws} therefore specializes to
        \[
            \alpha_v\circ F_1(a)
            =
            G_1(a)\circ\alpha_u,
        \]
        which is precisely internal naturality.

        \item \textbf{Reverse the construction.} Conversely, an internal natural transformation \(\alpha:F\Rightarrow G\) defines the span map
        \[
            \bar\alpha=\langle 1_{A_0},\alpha\rangle:
            A_0\to A_0\times_{F_0,B_0,s_B}B_1,
        \]
        and the naturality equation is exactly the Mealy simulation law. Hence the two notions are equivalent.
    \end{enumerate}
\end{proof}

\section{The bicategory of Mealy morphisms}
\label{sec:bicategory-of-mealy-morphisms}

Mealy morphisms compose by composition of oriented monad morphisms in spans. We first give the explicit formula.

Let
\[
    (P,\delta_P,\omega_P):\mathbb A\rightsquigarrow\mathbb B
\]
and
\[
    (Q,\delta_Q,\omega_Q):\mathbb B\rightsquigarrow\mathbb C
\]
be Mealy morphisms, with underlying spans
\[
    A_0\xleftarrow{p_P}P\xrightarrow{q_P}B_0,
    \qquad
    B_0\xleftarrow{p_Q}Q\xrightarrow{q_Q}C_0.
\]
Their composite carrier is the composite span
\[
    Q\circ P:
    A_0
    \xleftarrow{p_P\pi_P}
    P\times_{q_P,B_0,p_Q}Q
    \xrightarrow{q_Q\pi_Q}
    C_0.
\]
The defining pullback is
\[
    \begin{tikzcd}
            P\times_{q_P,B_0,p_Q}Q
                \arrow[r, "\pi_Q"]
                \arrow[d, "\pi_P"']
                \arrow[dr, phantom, "\lrcorner", very near start]
        &
            Q   
                \arrow[d, "p_Q"]
        \\
            P 
                \arrow[r, "q_P"']
        &
            B_0.
    \end{tikzcd}
\]
Operationally, a composite state is a pair
\[
    (x,y)
\]
satisfying
\[
    q_P(x)=p_Q(y).
\]
This is the pipeline diagram
\[
    \begin{tikzcd}[column sep=large]
            A_0
        &
            P 
                \arrow[l, "p_P"'] 
                \arrow[r, "q_P"]
        &
            B_0
        &
            Q 
                \arrow[l, "p_Q"'] 
                \arrow[r, "q_Q"]
        &
            C_0.
    \end{tikzcd}
\]

Given an input arrow
\[
    a:u\to v
\]
of \(\mathbb A\) with
\[
    p_P(x)=v,
\]
the first machine updates \(x\) to
\[
    \delta_P(a,x)
\]
and emits a \(\mathbb B\)-arrow
\[
    \omega_P(a,x):q_P(\delta_P(a,x))\to q_P(x).
\]
Since
\[
    q_P(x)=p_Q(y),
\]
this emitted arrow is admissible input for the second machine at \(y\). The second machine updates \(y\) to
\[
    \delta_Q(\omega_P(a,x),y)
\]
and emits
\[
    \omega_Q(\omega_P(a,x),y).
\]
Thus the composite transition and output are
\[
    \delta_{QP}(a,(x,y))
    =
    \bigl(
    \delta_P(a,x),
    \delta_Q(\omega_P(a,x),y)
    \bigr),
\]
and
\[
    \omega_{QP}(a,(x,y))
    =
    \omega_Q(\omega_P(a,x),y).
\]
The execution picture is
\[
    \begin{tikzcd}[column sep=large, row sep=large]
            x
                \arrow[r, "{a/\omega_P(a,x)}"]
        &
            \delta_P(a,x)
        \\
            y
                \arrow[r, "{\omega_P(a,x)/\omega_Q(\omega_P(a,x),y)}"']
        &
            \delta_Q(\omega_P(a,x),y).
    \end{tikzcd}
\]

\begin{proposition}[Pipeline composition]
\label{prop:pipeline-composition-mealy}
    The formulas above define a Mealy morphism
    \[
        Q\circ P:\mathbb A\rightsquigarrow\mathbb C.
    \]
\end{proposition}
\begin{proof}
    \leavevmode
    \begin{enumerate}
        \item \textbf{Check that the transition lands in the pullback carrier.} We must show that the two updated states have matching middle interface. By the typing equations for \(Q\),
        \[
            p_Q(\delta_Q(\omega_P(a,x),y))
            =
            s_B(\omega_P(a,x)).
        \]
        By the typing equations for \(P\),
        \[
            s_B(\omega_P(a,x))
            =
            q_P(\delta_P(a,x)).
        \]
        Hence
        \[
            p_Q(\delta_Q(\omega_P(a,x),y))
            =
            q_P(\delta_P(a,x)),
        \]
        so
        \[
            \bigl(
            \delta_P(a,x),
            \delta_Q(\omega_P(a,x),y)
            \bigr)
        \]
        lies in
        \[
            P\times_{B_0}Q.
        \]

        \item \textbf{Check the output typing.} The output of the composite is
        \[
            \omega_Q(\omega_P(a,x),y).
        \]
        Since \(Q\) is a Mealy morphism,
        \[
            s_C(\omega_Q(\omega_P(a,x),y))
            =
            q_Q(\delta_Q(\omega_P(a,x),y)),
        \]
        and
        \[
            t_C(\omega_Q(\omega_P(a,x),y))
            =
            q_Q(y).
        \]
        Thus the output arrow has the required type
        \[
            q_Q(\delta_Q(\omega_P(a,x),y))\to q_Q(y).
        \]

        \item \textbf{Verify the unit law.} Let \(x\in P\), \(y\in Q\), and suppose \(q_P(x)=p_Q(y)\). For the identity input arrow \(e_A(p_Px)\), the unit law for \(P\) gives
        \[
            \delta_P(e_A(p_Px),x)=x,
            \qquad
            \omega_P(e_A(p_Px),x)=e_B(q_Px).
        \]
        The unit law for \(Q\) then gives
        \[
            \delta_Q(e_B(q_Px),y)=y,
            \qquad
            \omega_Q(e_B(q_Px),y)=e_C(q_Qy).
        \]
        Hence
        \[
            \delta_{QP}(e_A(p_Px),(x,y))=(x,y),
        \]
        and
        \[
            \omega_{QP}(e_A(p_Px),(x,y))=e_C(q_Qy).
        \]

        \item \textbf{Prepare the multiplication calculation.} Let
        \[
            a:u\to v,
            \qquad
            b:v\to w,
            \qquad
            p_P(x)=w.
        \]
        Put
        \[
            \beta_b:=\omega_P(b,x),
            \qquad
            \beta_a:=\omega_P(a,\delta_P(b,x)).
        \]
        Then
        \[
            \beta_a:q_P(\delta_P(a,\delta_P(b,x)))\to q_P(\delta_P(b,x)),
        \]
        and
        \[
            \beta_b:q_P(\delta_P(b,x))\to q_P(x).
        \]
        By the Mealy multiplication law for \(P\),
        \[
            \omega_P(b\circ a,x)=\beta_b\circ\beta_a.
        \]

        \item \textbf{Verify the state part of multiplication.} The state part of executing the composite input \(b\circ a\) is
        \[
            \delta_{QP}(b\circ a,(x,y))
            =
            \left(
                \delta_P(b\circ a,x),
                \delta_Q(\omega_P(b\circ a,x),y)
            \right).
        \]
        By the Mealy multiplication law for \(P\),
        \[
            \delta_P(b\circ a,x)=\delta_P(a,\delta_P(b,x)).
        \]
        By the Mealy multiplication law for \(Q\), applied to the composite
        \[
            \beta_b\circ\beta_a,
        \]
        we have
        \[
            \delta_Q(\beta_b\circ\beta_a,y)
            =
            \delta_Q(\beta_a,\delta_Q(\beta_b,y)).
        \]
        This is exactly the \(Q\)-state obtained by first executing \(b\), then executing \(a\), in the composite machine.

        \item \textbf{Verify the output part of multiplication.} The output for the composite input is
        \[
            \omega_{QP}(b\circ a,(x,y))
            =
            \omega_Q(\omega_P(b\circ a,x),y).
        \]
        Since
        \[
            \omega_P(b\circ a,x)=\beta_b\circ\beta_a,
        \]
        the Mealy multiplication law for \(Q\) gives
        \[
            \omega_Q(\beta_b\circ\beta_a,y)
            =
            \omega_Q(\beta_b,y)\circ
            \omega_Q(\beta_a,\delta_Q(\beta_b,y)).
        \]
        This is precisely the composite of the outputs produced by first executing \(b\), then executing \(a\), in the pipeline. Hence the multiplication law holds.
    \end{enumerate}
\end{proof}

The identity Mealy morphism on \(\mathbb A\) has underlying identity span
\[
    A_0\xleftarrow{1_{A_0}}A_0\xrightarrow{1_{A_0}}A_0.
\]
At a state \(v\in A_0\), an arrow
\[
    a:u\to v
\]
is processed by
\[
    v
    \xrightarrow{\;a/a\;}
    u.
\]
Thus the identity Mealy morphism has
\[
    \delta(a,v)=s_A(a),
    \qquad
    \omega(a,v)=a.
\]
The corresponding diagram is
\[
    \begin{tikzcd}[column sep=large]
            u 
                \arrow[r, "a"] 
        & 
            v
    \end{tikzcd}
    \qquad\leadsto\qquad
    \begin{tikzcd}[column sep=large]
            v 
                \arrow[r, "{a/a}"] 
        & 
            u.
    \end{tikzcd}
\]

Recall that the displayed simulation direction is opposite to the underlying Kleisli algebra-morphism direction: a simulation
\[
    P\Rightarrow Q
\]
has underlying state map
\[
    Q\to P.
\]

\begin{proposition}[Vertical composition of simulations]
\label{prop:mealy-simulations-vertical-composition}
    Mealy simulations compose vertically. If
    \[
        (\theta,\alpha):P\Rightarrow Q
    \]
    and
    \[
        (\kappa,\beta):Q\Rightarrow R
    \]
    are simulations between Mealy morphisms
    \[
        P,Q,R:\mathbb A\rightsquigarrow\mathbb B,
    \]
    where
    \[
        \theta:Q\to P,
        \qquad
        \kappa:R\to Q,
    \]
    then their vertical composite
    \[
        (\kappa,\beta)\circ(\theta,\alpha):P\Rightarrow R
    \]
    has state map
    \[
        \theta\kappa:R\to P
    \]
    and comparison component
    \[
        \gamma_z=\beta_z\circ\alpha_{\kappa z}.
    \]
\end{proposition}
\begin{proof}
    \leavevmode
    \begin{enumerate}
        \item \textbf{Use composition of Kleisli algebra morphisms.}
        A simulation
        \[
            (\theta,\alpha):P\Rightarrow Q
        \]
        is, by definition, a morphism of
        \(\overline R_A\)-algebras
        \[
            (Q,\Psi)\to(P,\Phi)
        \]
        in \(\operatorname{Kl}(T_B)\). Similarly,
        \[
            (\kappa,\beta):Q\Rightarrow R
        \]
        is a morphism of algebras
        \[
            (R,\Xi)\to(Q,\Psi).
        \]
        Their composite in the algebra category is therefore a morphism
        \[
            (R,\Xi)\to(P,\Phi),
        \]
        which, under the displayed simulation convention, is a simulation
        \[
            P\Rightarrow R.
        \]

        \item \textbf{Unpack the Kleisli composite.}
        The underlying Kleisli arrows are represented by
        \[
            R\xrightarrow{(\kappa,\beta)}T_B(Q)
        \]
        and
        \[
            Q\xrightarrow{(\theta,\alpha)}T_B(P).
        \]
        Their Kleisli composite has state map
        \[
            \theta\kappa:R\to P.
        \]
        At a state \(z\in R\), its comparison component is the composite
        \[
            \begin{tikzcd}[column sep=large]
                    q_P(\theta\kappa z)
                        \arrow[r, "\alpha_{\kappa z}"]
                &
                    q_Q(\kappa z)
                        \arrow[r, "\beta_z"]
                &
                    q_R(z),
            \end{tikzcd}
        \]
        namely
        \[
            \gamma_z=\beta_z\circ\alpha_{\kappa z}.
        \]

        \item \textbf{Conclude.}
        Morphisms of \(\overline R_A\)-algebras are closed under composition.
        Hence the pair
        \[
            (\theta\kappa,\gamma)
        \]
        is again a Mealy simulation, and it is the vertical composite
        \[
            (\kappa,\beta)\circ(\theta,\alpha).
        \]
    \end{enumerate}
\end{proof}

\begin{proposition}[Horizontal composition of simulations]
\label{prop:mealy-simulations-horizontal-composition}
    Let
    \[
        P,P':\mathbb A\rightsquigarrow\mathbb B,
        \qquad
        Q,Q':\mathbb B\rightsquigarrow\mathbb C
    \]
    be Mealy morphisms. Suppose
    \[
        (\theta,\alpha):P\Rightarrow P',
        \qquad
        (\kappa,\beta):Q\Rightarrow Q'
    \]
    are Mealy simulations, so that
    \[
        \theta:P'\to P,
        \qquad
        \kappa:Q'\to Q.
    \]
    Their horizontal composite
    \[
        (\kappa,\beta)\ast(\theta,\alpha):
        Q\circ P\Rightarrow Q'\circ P'
    \]
    has state map
    \[
        P'\times_{B_0}Q'
        \longrightarrow
        P\times_{B_0}Q
    \]
    given by
    \[
        (x',y')
        \longmapsto
        \bigl(
            \theta x',
            \delta_Q(\alpha_{x'},\kappa y')
        \bigr),
    \]
    and comparison component
    \[
        \beta_{y'}\circ
        \omega_Q(\alpha_{x'},\kappa y').
    \]
\end{proposition}
\begin{proof}
    \leavevmode
    \begin{enumerate}
        \item \textbf{Check that the state map is defined.} Let \((x',y')\in P'\times_{B_0}Q'\), so
        \[
            q_{P'}(x')=p_{Q'}(y').
        \]
        Since \(p_Q\kappa=p_{Q'}\), the comparison arrow
        \[
            \alpha_{x'}:q_P(\theta x')\to q_{P'}(x')
        \]
        is admissible input for \(Q\) at \(\kappa y'\). Hence \(\delta_Q(\alpha_{x'},\kappa y')\) is defined, and its input interface is
        \[
            p_Q\delta_Q(\alpha_{x'},\kappa y')
            =s_B\alpha_{x'}
            =q_P\theta x'.
        \]
        Thus
        \[
            \bigl(\theta x',\delta_Q(\alpha_{x'},\kappa y')\bigr)
        \]
        lies in \(P\times_{B_0}Q\).

        \item \textbf{Check the comparison arrow.} The output of \(Q\) while processing \(\alpha_{x'}\) at \(\kappa y'\) is
        \[
            \omega_Q(\alpha_{x'},\kappa y'):
            q_Q\delta_Q(\alpha_{x'},\kappa y')
            \to
            q_Q\kappa y'.
        \]
        Composing with
        \[
            \beta_{y'}:q_Q\kappa y'\to q_{Q'}y'
        \]
        gives the required comparison arrow
        \[
            \beta_{y'}\circ\omega_Q(\alpha_{x'},\kappa y'):
            q_Q\delta_Q(\alpha_{x'},\kappa y')
            \to
            q_{Q'}y'.
        \]

        \item \textbf{Verify the simulation law.} Let \(a:u\to v\) be an input arrow for \(\mathbb A\), and let \(p_{P'}(x')=v\). Put
        \[
            x_1':=\delta_{P'}(a,x'),
            \qquad
            \varphi:=\omega_{P'}(a,x'),
            \qquad
            y_1':=\delta_{Q'}(\varphi,y').
        \]
        The state equation for \((\theta,\alpha)\) gives
        \[
            \theta x_1'=\delta_P(a,\theta x').
        \]
        The output equation for \((\theta,\alpha)\) gives
        \[
            \alpha_{x'}\circ\omega_P(a,\theta x')
            =
            \varphi\circ\alpha_{x_1'}.
        \]
        Applying the Mealy multiplication law for \(Q\) to this equality and then the state equation for \((\kappa,\beta)\) gives
        \[
            \delta_Q\bigl(
                \omega_P(a,\theta x'),
                \delta_Q(\alpha_{x'},\kappa y')
            \bigr)
            =
            \delta_Q(\alpha_{x_1'},\kappa y_1').
        \]
        Hence the state part of the simulation law holds.

        For the output part, set
        \[
            \zeta_{(x',y')}
            :=
            \beta_{y'}\circ\omega_Q(\alpha_{x'},\kappa y').
        \]
        The required equation is obtained by pasting the output square for \((\theta,\alpha)\), the Mealy multiplication law for \(Q\), and the output square for \((\kappa,\beta)\). Explicitly, the left-hand output composite reduces to
        \[
        \begin{aligned}
            &\beta_{y'}\circ
            \omega_Q(
                \alpha_{x'}\circ\omega_P(a,\theta x'),
                \kappa y'
            )
            \\
            &\qquad=
            \beta_{y'}\circ
            \omega_Q(
                \varphi\circ\alpha_{x_1'},
                \kappa y'
            )
            \\
            &\qquad=
            \omega_{Q'}(\varphi,y')\circ
            \beta_{y_1'}\circ
            \omega_Q(\alpha_{x_1'},\kappa y_1'),
        \end{aligned}
        \]
        which is
        \[
            \omega_{Q'\circ P'}(a,(x',y'))\circ\zeta_{(x_1',y_1')}.
        \]
        This is exactly the simulation output equation.
    \end{enumerate}
\end{proof}

\begin{proposition}[Functoriality of horizontal composition of simulations]
\label{prop:mealy-horizontal-composition-functorial}
    Let
    \[
        P,P',P'':\mathbb A\rightsquigarrow\mathbb B,
        \qquad
        Q,Q',Q'':\mathbb B\rightsquigarrow\mathbb C
    \]
    be Mealy morphisms. Suppose we are given vertically composable Mealy simulations
    \[
        (\theta,\alpha):P\Rightarrow P',
        \qquad
        (\theta',\alpha'):P'\Rightarrow P'',
    \]
    and
    \[
        (\kappa,\beta):Q\Rightarrow Q',
        \qquad
        (\kappa',\beta'):Q'\Rightarrow Q''.
    \]
    Thus
    \[
        \theta:P'\to P,
        \qquad
        \theta':P''\to P',
        \qquad
        \kappa:Q'\to Q,
        \qquad
        \kappa':Q''\to Q'
    \]
    are the underlying state maps. Horizontal composition of simulations preserves identities and vertical composites. Hence, for every triple of internal categories
    \[
        \mathbb A,\mathbb B,\mathbb C,
    \]
    pipeline composition defines a functor
    \[
        \mathbf{Mly}(\mathcal E)(\mathbb B,\mathbb C)
        \times
        \mathbf{Mly}(\mathcal E)(\mathbb A,\mathbb B)
        \longrightarrow
        \mathbf{Mly}(\mathcal E)(\mathbb A,\mathbb C).
    \]
\end{proposition}
\begin{proof}
    \leavevmode
    \begin{enumerate}
        \item \textbf{Identities.}
        If both simulations are identity simulations, then
        \[
            \theta=1_P,
            \quad
            \alpha_x=e_B(q_Px),
            \qquad
            \kappa=1_Q,
            \quad
            \beta_y=e_C(q_Qy).
        \]
        The formula of
        Proposition~\ref{prop:mealy-simulations-horizontal-composition}
        sends \((x,y)\in P\times_{B_0}Q\) to
        \[
            \bigl(x,\delta_Q(e_B(q_Px),y)\bigr)
            =
            (x,y)
        \]
        by the Mealy unit law for \(Q\). Its comparison component is
        \[
            e_C(q_Qy)\circ
            \omega_Q(e_B(q_Px),y)
            =
            e_C(q_Qy).
        \]
        Hence horizontal composition preserves identity simulations.

        \item \textbf{Vertical composites and their horizontal composite.}
        The vertical composites
        \[
            (\theta',\alpha')\circ(\theta,\alpha):
            P\Rightarrow P''
        \]
        and
        \[
            (\kappa',\beta')\circ(\kappa,\beta):
            Q\Rightarrow Q''
        \]
        have state maps
        \[
            \theta\theta':P''\to P,
            \qquad
            \kappa\kappa':Q''\to Q,
        \]
        and comparison components
        \[
            \gamma_{x''}
            :=
            \alpha'_{x''}\circ
            \alpha_{\theta'x''},
        \]
        \[
            \eta_{y''}
            :=
            \beta'_{y''}\circ
            \beta_{\kappa'y''}.
        \]

        Horizontally composing these vertical composites gives the state map
        \[
            (x'',y'')
            \longmapsto
            \left(
                \theta\theta'x'',
                \delta_Q(\gamma_{x''},\kappa\kappa'y'')
            \right)
        \]
        and comparison component
        \[
            \eta_{y''}\circ
            \omega_Q(\gamma_{x''},\kappa\kappa'y'').
        \]

        \item \textbf{Horizontally compose first.}
        Horizontally composing
        \[
            (\theta,\alpha):P\Rightarrow P'
        \]
        with
        \[
            (\kappa,\beta):Q\Rightarrow Q'
        \]
        gives a simulation
        \[
            Q\circ P\Rightarrow Q'\circ P'
        \]
        whose state map is
        \[
            (x',y')
            \longmapsto
            \left(
                \theta x',
                \delta_Q(\alpha_{x'},\kappa y')
            \right)
        \]
        and whose comparison component is
        \[
            \beta_{y'}\circ
            \omega_Q(\alpha_{x'},\kappa y').
        \]

        Similarly, horizontally composing
        \[
            (\theta',\alpha'):P'\Rightarrow P''
        \]
        with
        \[
            (\kappa',\beta'):Q'\Rightarrow Q''
        \]
        gives a simulation
        \[
            Q'\circ P'\Rightarrow Q''\circ P''
        \]
        whose state map is
        \[
            (x'',y'')
            \longmapsto
            \left(
                \theta'x'',
                \delta_{Q'}(\alpha'_{x''},\kappa'y'')
            \right)
        \]
        and whose comparison component is
        \[
            \beta'_{y''}\circ
            \omega_{Q'}(\alpha'_{x''},\kappa'y'').
        \]

        Put
        \[
            \overline y
            :=
            \delta_{Q'}(\alpha'_{x''},\kappa'y'').
        \]

        \item \textbf{Compare the state maps.}
        Vertically composing the two horizontally composed simulations gives
        the state
        \[
            \left(
                \theta\theta'x'',
                \delta_Q(
                    \alpha_{\theta'x''},
                    \kappa\overline y
                )
            \right).
        \]
        The transition equation for the simulation
        \((\kappa,\beta):Q\Rightarrow Q'\), applied to the input arrow
        \(\alpha'_{x''}\), gives
        \[
            \kappa\overline y
            =
            \delta_Q(
                \alpha'_{x''},
                \kappa\kappa'y''
            ).
        \]
        Therefore the second state component is
        \[
            \delta_Q\!\left(
                \alpha_{\theta'x''},
                \delta_Q(
                    \alpha'_{x''},
                    \kappa\kappa'y''
                )
            \right).
        \]
        By the Mealy multiplication law for \(Q\),
        \[
        \begin{aligned}
            &
            \delta_Q\!\left(
                \alpha_{\theta'x''},
                \delta_Q(
                    \alpha'_{x''},
                    \kappa\kappa'y''
                )
            \right)
            \\
            &\qquad=
            \delta_Q\!\left(
                \alpha'_{x''}\circ\alpha_{\theta'x''},
                \kappa\kappa'y''
            \right)
            \\
            &\qquad=
            \delta_Q(
                \gamma_{x''},
                \kappa\kappa'y''
            ).
        \end{aligned}
        \]
        Thus the two constructions have the same state map.

        \item \textbf{Compare the output components.}
        The output-comparison component obtained by horizontally composing
        first and then vertically composing is
        \[
        \begin{aligned}
            &
            \beta'_{y''}
            \circ
            \omega_{Q'}(
                \alpha'_{x''},
                \kappa'y''
            )
            \circ
            \beta_{\overline y}
            \circ
            \omega_Q(
                \alpha_{\theta'x''},
                \kappa\overline y
            ).
        \end{aligned}
        \]
        The output equation for the simulation
        \((\kappa,\beta):Q\Rightarrow Q'\), applied to
        \(\alpha'_{x''}\) at the state \(\kappa'y''\), is
        \[
            \beta_{\kappa'y''}
            \circ
            \omega_Q(
                \alpha'_{x''},
                \kappa\kappa'y''
            )
            =
            \omega_{Q'}(
                \alpha'_{x''},
                \kappa'y''
            )
            \circ
            \beta_{\overline y}.
        \]
        Using this equation and the previously established identity
        \[
            \kappa\overline y
            =
            \delta_Q(
                \alpha'_{x''},
                \kappa\kappa'y''
            ),
        \]
        the output component becomes
        \[
        \begin{aligned}
            &
            \beta'_{y''}
            \circ
            \beta_{\kappa'y''}
            \circ
            \omega_Q(
                \alpha'_{x''},
                \kappa\kappa'y''
            )
            \circ
            \omega_Q\!\left(
                \alpha_{\theta'x''},
                \delta_Q(
                    \alpha'_{x''},
                    \kappa\kappa'y''
                )
            \right)
            \\
            &\qquad=
            \eta_{y''}
            \circ
            \omega_Q\!\left(
                \alpha'_{x''}\circ\alpha_{\theta'x''},
                \kappa\kappa'y''
            \right)
            \\
            &\qquad=
            \eta_{y''}
            \circ
            \omega_Q(
                \gamma_{x''},
                \kappa\kappa'y''
            ),
        \end{aligned}
        \]
        where the first equality uses the Mealy output multiplication law for
        \(Q\).

        This is exactly the comparison component obtained by first taking the
        vertical composites and then horizontally composing them.

        \item \textbf{Conclusion.}
        Horizontal composition preserves identity simulations and vertical
        composites. Hence, for every triple of internal categories
        \[
            \mathbb A,\mathbb B,\mathbb C,
        \]
        pipeline composition defines a functor
        \[
            \mathbf{Mly}(\mathcal E)(\mathbb B,\mathbb C)
            \times
            \mathbf{Mly}(\mathcal E)(\mathbb A,\mathbb B)
            \longrightarrow
            \mathbf{Mly}(\mathcal E)(\mathbb A,\mathbb C).
        \]
    \end{enumerate}
\end{proof}

\begin{lemma}[Span coherence induces strict Mealy coherence]
\label{lem:span-coherence-strict-mealy}
    The canonical associativity and unit isomorphisms for pullback composition
    determine invertible strict Mealy simulations.

    More precisely, because the displayed direction of a Mealy simulation is
    opposite to the direction of its underlying state map, the state map of
    each Mealy coherence simulation is the inverse of the apex map of the
    corresponding coherence 2-cell in \(\mathbf{Span}(\mathcal E)\).
\end{lemma}
\begin{proof}
    \leavevmode
    \begin{enumerate}
        \item \textbf{Associators.}
        Let
        \[
            P:\mathbb A\rightsquigarrow\mathbb B,
            \qquad
            Q:\mathbb B\rightsquigarrow\mathbb C,
            \qquad
            R:\mathbb C\rightsquigarrow\mathbb D
        \]
        be composable Mealy morphisms.

        The carrier of
        \[
            (R\circ Q)\circ P
        \]
        is
        \[
            P\times_{B_0}(Q\times_{C_0}R),
        \]
        whereas the carrier of
        \[
            R\circ(Q\circ P)
        \]
        is
        \[
            (P\times_{B_0}Q)\times_{C_0}R.
        \]
        Both are iterated-pullback presentations of the object of triples
        \[
            (x,y,z)
        \]
        satisfying
        \[
            q_P(x)=p_Q(y),
            \qquad
            q_Q(y)=p_R(z).
        \]

        The ordinary span associator
        \[
            (R\circ Q)\circ P
            \Longrightarrow
            R\circ(Q\circ P)
        \]
        has apex map
        \[
            P\times_{B_0}(Q\times_{C_0}R)
            \longrightarrow
            (P\times_{B_0}Q)\times_{C_0}R,
        \]
        characterized in generalized-element notation by
        \[
            (x,(y,z))
            \longmapsto
            ((x,y),z).
        \]

        We define the Mealy associator
        \[
            a_{R,Q,P}:
            (R\circ Q)\circ P
            \Rightarrow
            R\circ(Q\circ P)
        \]
        using the inverse pullback isomorphism as its state map:
        \[
            \widehat a_{P,Q,R}:
            (P\times_{B_0}Q)\times_{C_0}R
            \xrightarrow{\cong}
            P\times_{B_0}(Q\times_{C_0}R),
        \]
        \[
            \widehat a_{P,Q,R}((x,y),z)
            =
            (x,(y,z)).
        \]
        This direction is forced by the convention that a simulation from a
        source Mealy morphism to a target Mealy morphism has a state map from
        the target carrier to the source carrier.

        Under the canonical identification of both carriers with the same
        triple-pullback object, the two parenthesized pipelines execute \(P\),
        then \(Q\), then \(R\). They therefore have the same transition map
        and emit the same output arrow. Consequently
        \(\widehat a_{P,Q,R}\), equipped with the identity output comparison,
        is a strict Mealy simulation.

        The ordinary span-associator apex map gives the state map of the
        inverse Mealy simulation
        \[
            R\circ(Q\circ P)
            \Rightarrow
            (R\circ Q)\circ P.
        \]
        Hence the Mealy associator is invertible.

        \item \textbf{Left unitors.}
        Let
        \[
            P:\mathbb A\rightsquigarrow\mathbb B.
        \]
        The carrier of
        \[
            1_{\mathbb B}\circ P
        \]
        is
        \[
            P\times_{q_P,B_0,1_{B_0}}B_0.
        \]
        The ordinary span left unitor
        \[
            1_{B_0}\circ P\Longrightarrow P
        \]
        has apex map the canonical projection
        \[
            P\times_{q_P,B_0,1_{B_0}}B_0\to P.
        \]

        The Mealy left unitor is displayed in the same direction,
        \[
            \ell_P:
            1_{\mathbb B}\circ P
            \Rightarrow
            P,
        \]
        but its state map must run from the target carrier to the source
        carrier. We therefore use the inverse pullback isomorphism
        \[
            \widehat\ell_P:
            P\longrightarrow
            P\times_{q_P,B_0,1_{B_0}}B_0,
            \qquad
            x\longmapsto(x,q_Px).
        \]

        After \(P\) executes an input arrow \(a\) at \(x\), it updates to
        \(\delta_P(a,x)\) and emits
        \[
            \omega_P(a,x):
            q_P(\delta_P(a,x))\to q_P(x).
        \]
        The identity Mealy morphism on \(\mathbb B\) processes this emitted
        arrow at the object-state \(q_P(x)\), updates that state to
        \(q_P(\delta_P(a,x))\), and emits the same arrow. Thus the updated
        composite state is
        \[
            \bigl(
                \delta_P(a,x),
                q_P(\delta_P(a,x))
            \bigr)
            =
            \widehat\ell_P(\delta_P(a,x)),
        \]
        and the composite output is \(\omega_P(a,x)\).

        Hence \(\widehat\ell_P\), equipped with the identity output comparison,
        is a strict Mealy simulation. The canonical projection
        \[
            P\times_{q_P,B_0,1_{B_0}}B_0\to P
        \]
        gives the state map of its inverse strict simulation
        \[
            P\Rightarrow1_{\mathbb B}\circ P.
        \]

        \item \textbf{Right unitors.}
        The carrier of
        \[
            P\circ1_{\mathbb A}
        \]
        is
        \[
            A_0\times_{1_{A_0},A_0,p_P}P.
        \]
        The ordinary span right unitor
        \[
            P\circ1_{A_0}\Longrightarrow P
        \]
        has apex map the canonical projection
        \[
            A_0\times_{1_{A_0},A_0,p_P}P\to P.
        \]

        The Mealy right unitor
        \[
            r_P:
            P\circ1_{\mathbb A}
            \Rightarrow
            P
        \]
        therefore uses the inverse pullback isomorphism as its state map:
        \[
            \widehat r_P:
            P\longrightarrow
            A_0\times_{1_{A_0},A_0,p_P}P,
            \qquad
            x\longmapsto(p_Px,x).
        \]

        Let
        \[
            a:u\to v
        \]
        be an input arrow with \(p_P(x)=v\). The identity Mealy morphism on
        \(\mathbb A\) updates its object-state from \(v\) to \(u\) and emits
        \(a\). The machine \(P\) then processes \(a\) at \(x\), updating to
        \(\delta_P(a,x)\) and emitting \(\omega_P(a,x)\). Since
        \[
            p_P(\delta_P(a,x))=u,
        \]
        the updated composite state is
        \[
            (u,\delta_P(a,x))
            =
            \widehat r_P(\delta_P(a,x)).
        \]
        Thus \(\widehat r_P\) strictly preserves both transition and output.

        Hence \(\widehat r_P\), equipped with the identity output comparison,
        is a strict Mealy simulation. The canonical projection
        \[
            A_0\times_{1_{A_0},A_0,p_P}P\to P
        \]
        gives the state map of its inverse strict simulation
        \[
            P\Rightarrow P\circ1_{\mathbb A}.
        \]

        \item \textbf{Invertibility.}
        In each case, the state maps of the two displayed simulations are
        inverse pullback isomorphisms, and every output-comparison component is
        an identity arrow. Hence the associators and unitors are invertible
        strict Mealy simulations.
    \end{enumerate}
\end{proof}

\begin{lemma}[Naturality of pipeline coherence]
\label{lem:pipeline-coherence-naturality}
    The associativity and unit simulations of
    Lemma~\ref{lem:span-coherence-strict-mealy} are natural with respect to
    Mealy simulations.
\end{lemma}
\begin{proof}
    \leavevmode
    \begin{enumerate}
        \item \textbf{Set up the associator naturality calculation.}
        Let
        \[
            (\theta,\alpha):P\Rightarrow P',
            \qquad
            (\kappa,\beta):Q\Rightarrow Q',
            \qquad
            (\ell,\chi):R\Rightarrow R'
        \]
        be simulations between composable Mealy morphisms
        \[
            P,P':\mathbb A\rightsquigarrow\mathbb B,
        \]
        \[
            Q,Q':\mathbb B\rightsquigarrow\mathbb C,
        \]
        and
        \[
            R,R':\mathbb C\rightsquigarrow\mathbb D.
        \]
        Thus the state maps have directions
        \[
            \theta:P'\to P,
            \qquad
            \kappa:Q'\to Q,
            \qquad
            \ell:R'\to R.
        \]

        Let
        \[
            (x',y',z')
        \]
        be a state of the target triple pipeline. Put
        \[
            \gamma_{x',y'}
            :=
            \omega_Q(\alpha_{x'},\kappa y')
        \]
        and
        \[
            \zeta_{x',y'}
            :=
            \beta_{y'}\circ\gamma_{x',y'}.
        \]
        Here \(\gamma_{x',y'}\) and \(\beta_{y'}\) are composable arrows in
        \(\mathbb C\).

        \item \textbf{Compare the state components for the associator.}
        After identifying the two parenthesized carrier objects with the same
        triple pullback, one route around the associator naturality square
        gives the state
        \[
            \left(
                \theta x',
                \delta_Q(\alpha_{x'},\kappa y'),
                \delta_R(\zeta_{x',y'},\ell z')
            \right).
        \]
        The other route gives
        \[
            \left(
                \theta x',
                \delta_Q(\alpha_{x'},\kappa y'),
                \delta_R\!\left(
                    \gamma_{x',y'},
                    \delta_R(\beta_{y'},\ell z')
                \right)
            \right).
        \]
        By the state multiplication law for \(R\),
        \[
            \delta_R(
                \beta_{y'}\circ\gamma_{x',y'},
                \ell z'
            )
            =
            \delta_R\!\left(
                \gamma_{x',y'},
                \delta_R(\beta_{y'},\ell z')
            \right).
        \]
        Since
        \[
            \zeta_{x',y'}
            =
            \beta_{y'}\circ\gamma_{x',y'},
        \]
        the two state components agree.

        \item \textbf{Compare the output components for the associator.}
        The first route has output-comparison component
        \[
            \chi_{z'}\circ
            \omega_R(\zeta_{x',y'},\ell z').
        \]
        The second route has output-comparison component
        \[
            \chi_{z'}\circ
            \omega_R(\beta_{y'},\ell z')
            \circ
            \omega_R\!\left(
                \gamma_{x',y'},
                \delta_R(\beta_{y'},\ell z')
            \right).
        \]
        By the output multiplication law for \(R\),
        \[
            \omega_R(
                \beta_{y'}\circ\gamma_{x',y'},
                \ell z'
            )
            =
            \omega_R(\beta_{y'},\ell z')
            \circ
            \omega_R\!\left(
                \gamma_{x',y'},
                \delta_R(\beta_{y'},\ell z')
            \right).
        \]
        Therefore the two output-comparison components agree. This proves
        naturality of the associator.

        \item \textbf{Naturality of the left unitor.}
        Let
        \[
            (\theta,\alpha):P\Rightarrow P'
        \]
        be a Mealy simulation between morphisms
        \[
            P,P':\mathbb A\rightsquigarrow\mathbb B.
        \]
        Thus
        \[
            \theta:P'\to P
        \]
        and
        \[
            \alpha_{x'}:
            q_P(\theta x')\to q_{P'}(x').
        \]

        Under the canonical pullback presentation, a state of
        \[
            1_{\mathbb B}\circ P'
        \]
        has the form
        \[
            (x',q_{P'}x').
        \]
        Horizontally composing \((\theta,\alpha)\) with the identity simulation
        on \(1_{\mathbb B}\) sends this state to
        \[
            \left(
                \theta x',
                \delta_{1_{\mathbb B}}
                    (\alpha_{x'},q_{P'}x')
            \right).
        \]
        The identity Mealy morphism on \(\mathbb B\) satisfies
        \[
            \delta_{1_{\mathbb B}}(\beta,t_B\beta)=s_B\beta,
            \qquad
            \omega_{1_{\mathbb B}}(\beta,t_B\beta)=\beta.
        \]
        Since
        \[
            s_B\alpha_{x'}=q_P(\theta x'),
            \qquad
            t_B\alpha_{x'}=q_{P'}(x'),
        \]
        the resulting state is
        \[
            (\theta x',q_P\theta x').
        \]
        Its output-comparison component is
        \[
            e_B(q_{P'}x')
            \circ
            \omega_{1_{\mathbb B}}
                (\alpha_{x'},q_{P'}x')
            =
            \alpha_{x'}.
        \]

        The other route around the left-unitor naturality square first applies
        the simulation \((\theta,\alpha)\) and then the left-unitor simulation
        for \(P\). It has the same state map
        \[
            x'\longmapsto(\theta x',q_P\theta x')
        \]
        and the same output-comparison component \(\alpha_{x'}\).
        Hence the left unitor is natural.

        \item \textbf{Naturality of the right unitor.}
        Under the canonical pullback presentation, a state of
        \[
            P'\circ1_{\mathbb A}
        \]
        has the form
        \[
            (p_{P'}x',x').
        \]
        Horizontally composing the identity simulation on \(1_{\mathbb A}\)
        with \((\theta,\alpha)\) sends this state to
        \[
            \left(
                p_{P'}x',
                \delta_P(e_A(p_{P'}x'),\theta x')
            \right).
        \]
        Since
        \[
            p_P\theta=p_{P'},
        \]
        the Mealy unit law for \(P\) gives
        \[
            \delta_P(e_A(p_{P'}x'),\theta x')
            =
            \theta x'.
        \]
        Thus the resulting state is
        \[
            (p_P\theta x',\theta x').
        \]
        Its output-comparison component is
        \[
            \alpha_{x'}
            \circ
            \omega_P(e_A(p_{P'}x'),\theta x')
            =
            \alpha_{x'}.
        \]

        The other route around the right-unitor naturality square first
        applies \((\theta,\alpha)\) and then the right-unitor simulation for
        \(P\). It has the same state map
        \[
            x'\longmapsto(p_P\theta x',\theta x')
        \]
        and the same output-comparison component \(\alpha_{x'}\).
        Hence the right unitor is natural.
    \end{enumerate}
\end{proof}

\begin{theorem}[The Mealy bicategory]
\label{thm:mealy-bicategory}
    If \(\mathcal E\) has pullbacks, then there is a bicategory
    \[
        \mathbf{Mly}(\mathcal E)
    \]
    whose objects are internal categories in \(\mathcal E\), whose 1-cells are
    Mealy morphisms, and whose 2-cells are Mealy simulations.
\end{theorem}
\begin{proof}
    \leavevmode
    \begin{enumerate}
        \item \textbf{Objects and hom-categories.}
        The objects are internal categories in \(\mathcal E\), equivalently
        monads in \(\mathbf{Span}(\mathcal E)\).

        For internal categories \(\mathbb A\) and \(\mathbb B\), the objects of
        \[
            \mathbf{Mly}(\mathcal E)(\mathbb A,\mathbb B)
        \]
        are Mealy morphisms
        \[
            \mathbb A\rightsquigarrow\mathbb B.
        \]
        By Theorem~\ref{thm:mealy-morphisms-kleisli-algebras}, these are
        algebras for the lifted monad
        \[
            \overline R_A
        \]
        in \(\operatorname{Kl}(T_B)\).

        A Mealy simulation
        \[
            (P,\Phi)\Rightarrow(Q,\Psi)
        \]
        is, by definition, an algebra morphism
        \[
            (Q,\Psi)\to(P,\Phi).
        \]
        Hence the hom-category is the opposite of the corresponding algebra
        category. Its identity morphisms and vertical composition are
        therefore inherited from the algebra category, with vertical
        composition given explicitly by
        Proposition~\ref{prop:mealy-simulations-vertical-composition}.

        \item \textbf{Horizontal composition.}
        Horizontal composition of 1-cells is pipeline composition. For
        \[
            P:\mathbb A\rightsquigarrow\mathbb B,
            \qquad
            Q:\mathbb B\rightsquigarrow\mathbb C,
        \]
        the composite carrier is
        \[
            A_0
            \xleftarrow{p_P\pi_P}
            P\times_{q_P,B_0,p_Q}Q
            \xrightarrow{q_Q\pi_Q}
            C_0,
        \]
        with transition and output maps
        \[
            \delta_{QP}(a,(x,y))
            =
            \bigl(
                \delta_P(a,x),
                \delta_Q(\omega_P(a,x),y)
            \bigr)
        \]
        and
        \[
            \omega_{QP}(a,(x,y))
            =
            \omega_Q(\omega_P(a,x),y).
        \]
        Proposition~\ref{prop:pipeline-composition-mealy} proves that these
        data satisfy the Mealy laws.

        Horizontal composition of 2-cells is the construction of
        Proposition~\ref{prop:mealy-simulations-horizontal-composition}.
        Proposition~\ref{prop:mealy-horizontal-composition-functorial} proves
        that this operation preserves identity simulations and vertical
        composites. Thus pipeline composition defines a functor on each pair
        of hom-categories.

        \item \textbf{Identity 1-cells.}
        The identity 1-cell on an internal category \(\mathbb A\) has carrier
        \[
            A_0
            \xleftarrow{1_{A_0}}
            A_0
            \xrightarrow{1_{A_0}}
            A_0.
        \]
        Its transition and output maps are
        \[
            \delta(a,v)=s_A(a),
            \qquad
            \omega(a,v)=a,
        \]
        where necessarily \(v=t_A(a)\).

        For the identity input at \(v\), one has
        \[
            \delta(e_A(v),v)
            =
            s_Ae_A(v)
            =
            v
        \]
        and
        \[
            \omega(e_A(v),v)
            =
            e_A(v),
        \]
        so the Mealy unit law holds.

        Now let
        \[
            a:u\to v,
            \qquad
            b:v\to w
        \]
        be composable arrows. The transition multiplication law is
        \[
        \begin{aligned}
            \delta(b\circ a,w)
            &=
            s_A(b\circ a)
            \\
            &=
            u
            \\
            &=
            s_A(a)
            \\
            &=
            \delta(a,s_A(b))
            \\
            &=
            \delta(a,\delta(b,w)).
        \end{aligned}
        \]
        The output multiplication law is
        \[
        \begin{aligned}
            \omega(b\circ a,w)
            &=
            b\circ a
            \\
            &=
            \omega(b,w)\circ
            \omega(a,\delta(b,w)).
        \end{aligned}
        \]
        Thus these formulas define the identity Mealy morphism on
        \(\mathbb A\).

        \item \textbf{Associators and unitors.}
        The associators and unitors are obtained from the canonical
        associativity and unit isomorphisms for pullback composition.
        Because the displayed direction of a Mealy simulation is opposite to
        the direction of its underlying state map, their state maps are the
        inverse apex maps of the corresponding span-coherence 2-cells.

        By Lemma~\ref{lem:span-coherence-strict-mealy}, these maps, equipped
        with identity output comparisons, are invertible strict Mealy
        simulations. By Lemma~\ref{lem:pipeline-coherence-naturality}, they are
        natural with respect to arbitrary Mealy simulations.

        \item \textbf{Interchange.}
        The interchange law is precisely the functoriality of horizontal
        composition established in
        Proposition~\ref{prop:mealy-horizontal-composition-functorial}.

        \item \textbf{Pentagon and triangle coherence.}
        The state maps of the associators and unitors are canonical
        isomorphisms between iterated pullback presentations. Their horizontal
        whiskerings by identity simulations again have the corresponding
        canonical pullback maps as their state components.

        Indeed, the comparison component of an identity simulation is an
        identity arrow. In the horizontal-composition formula of
        Proposition~\ref{prop:mealy-simulations-horizontal-composition}, the
        additional transition and output terms introduced by whiskering
        therefore reduce by the Mealy unit laws:
        \[
            \delta(e(-),x)=x,
            \qquad
            \omega(e(-),x)=e(-).
        \]
        Consequently, whiskering an associator or unitor coherence simulation
        by an identity simulation again produces a strict simulation with
        identity output-comparison component.

        Every route in the pentagon therefore determines a canonical map from
        the carrier of the final parenthesization to the carrier of the
        initial parenthesization. The maps obtained along the two routes have
        the same projections to the four original carrier objects, and hence
        are equal by the universal properties of the relevant iterated
        pullbacks.

        The same argument applies to the two routes in the triangle: both
        state maps have the same projections to the original carrier objects,
        and hence coincide.

        Every output-comparison component occurring in the pentagon and
        triangle is an identity arrow. Their vertical composites are therefore
        identity arrows as well. Hence the pentagon and triangle identities
        hold as equations of Mealy simulations.

        \item \textbf{Conclusion.}
        The preceding data provide hom-categories, functorial horizontal
        composition, identity 1-cells, natural invertible associators and
        unitors, and the required interchange and coherence equations. They
        therefore define a bicategory
        \[
            \mathbf{Mly}(\mathcal E)
        \]
        whose objects are internal categories, whose 1-cells are Mealy
        morphisms, and whose 2-cells are Mealy simulations.
    \end{enumerate}
\end{proof}

\begin{corollary}[The functorial sub-bicategory]
\label{cor:functorial-sub-bicategory}
    Let
    \[
        \mathbf{Mly}_{\mathrm{map}}(\mathcal E)
    \]
    denote the locally full sub-bicategory of
    \[
        \mathbf{Mly}(\mathcal E)
    \]
    having all internal categories as objects and the map-representable Mealy
    morphisms as 1-cells. Then
    \[
        \mathbf{Mly}_{\mathrm{map}}(\mathcal E)
    \]
    is biequivalent to the ordinary 2-category
    \[
        \mathbf{Cat}(\mathcal E)
    \]
    of internal categories, internal functors, and internal natural
    transformations.

    More precisely, after restricting to the canonical carriers
    \[
        F_{0!}
        =
        \left(
            A_0
            \xleftarrow{1_{A_0}}
            A_0
            \xrightarrow{F_0}
            B_0
        \right)
    \]
    and transporting pipeline composition along the canonical pullback
    isomorphisms, one obtains the usual strict 2-categorical composition of
    internal functors and internal natural transformations.
\end{corollary}
\begin{proof}
    \leavevmode
    \begin{enumerate}
        \item \textbf{Map-representable carriers form a sub-bicategory.}
        Identity carriers are map-representable, since the identity carrier on
        \(\mathbb A\) is
        \[
            A_0
            \xleftarrow{1_{A_0}}
            A_0
            \xrightarrow{1_{A_0}}
            A_0.
        \]

        It remains to check closure under pipeline composition. Let
        \[
            A_0\xleftarrow{p_P}P\xrightarrow{q_P}B_0,
            \qquad
            B_0\xleftarrow{p_Q}Q\xrightarrow{q_Q}C_0
        \]
        be map-representable carriers. Thus \(p_P\) and \(p_Q\) are
        isomorphisms. Their composite carrier has apex
        \[
            P\times_{q_P,B_0,p_Q}Q
        \]
        and left leg
        \[
            p_P\pi_P:
            P\times_{q_P,B_0,p_Q}Q\to A_0.
        \]
        In the defining pullback square
        \[
            \begin{tikzcd}
                    P\times_{q_P,B_0,p_Q}Q
                        \arrow[r, "\pi_Q"]
                        \arrow[d, "\pi_P"']
                        \arrow[dr, phantom, "\lrcorner", very near start]
                &
                    Q
                        \arrow[d, "p_Q"]
                \\
                    P
                        \arrow[r, "q_P"']
                &
                    B_0,
            \end{tikzcd}
        \]
        the projection \(\pi_P\) is a pullback of the isomorphism \(p_Q\), and
        is therefore an isomorphism. Hence \(p_P\pi_P\) is an isomorphism, so
        the composite carrier is map-representable.

        Since the sub-bicategory is locally full, it contains every Mealy
        simulation between its 1-cells. Therefore
        \[
            \mathbf{Mly}_{\mathrm{map}}(\mathcal E)
        \]
        is well-defined.

        \item \textbf{Canonical map carriers correspond to internal functors.}
        By Theorem~\ref{thm:representable-mealy-internal-functor}, for every
        morphism
        \[
            F_0:A_0\to B_0,
        \]
        Mealy structures on the canonical carrier
        \[
            F_{0!}
            =
            \left(
                A_0
                \xleftarrow{1_{A_0}}
                A_0
                \xrightarrow{F_0}
                B_0
            \right)
        \]
        are in bijective correspondence with internal functors
        \[
            F:\mathbb A\to\mathbb B
        \]
        having object map \(F_0\).

        \item \textbf{Every map-representable carrier is isomorphic to a
        canonical carrier in the Mealy hom-category.}
        Let
        \[
            (P,\Phi):\mathbb A\rightsquigarrow\mathbb B
        \]
        have map-representable carrier
        \[
            A_0\xleftarrow{p}P\xrightarrow{q}B_0.
        \]
        Then \(p\) is an isomorphism. Put
        \[
            F_0:=qp^{-1}:A_0\to B_0.
        \]
        By
        Corollary~\ref{cor:map-representable-mealy-internal-functor-up-to-iso},
        the carrier isomorphism
        \[
            p:
            \left(
                A_0\xleftarrow{p}P\xrightarrow{q}B_0
            \right)
            \overset{\cong}{\Longrightarrow}
            F_{0!}
        \]
        transports \(\Phi\) to the Mealy structure
        \[
            \Phi^F
            :=
            (Bp)\circ\Phi\circ(p^{-1}A)
        \]
        on \(F_{0!}\).

        We first construct a strict simulation
        \[
            (P,\Phi)\Rightarrow(F_{0!},\Phi^F).
        \]
        Its state map is
        \[
            p^{-1}:A_0\to P.
        \]
        This is a map of carrier spans because
        \[
            p\,p^{-1}=1_{A_0}
        \]
        and
        \[
            qp^{-1}=F_0.
        \]
        Its output-comparison component is therefore the identity arrow
        \[
            e_BF_0:A_0\to B_1.
        \]

        The strict compatibility equation for this state map is
        \[
            (Bp^{-1})\circ\Phi^F
            =
            \Phi\circ(p^{-1}A).
        \]
        Indeed,
        \[
        \begin{aligned}
            (Bp^{-1})\circ\Phi^F
            &=
            (Bp^{-1})\circ
            (Bp)\circ
            \Phi\circ
            (p^{-1}A)
            \\
            &=
            \Phi\circ(p^{-1}A).
        \end{aligned}
        \]
        Here and below the canonical associativity constraints are suppressed.

        Conversely, define a strict simulation
        \[
            (F_{0!},\Phi^F)\Rightarrow(P,\Phi)
        \]
        with state map
        \[
            p:P\to A_0.
        \]
        Its output-comparison component is the identity arrow
        \[
            e_Bq:P\to B_1.
        \]
        The required strict compatibility equation is
        \[
            (Bp)\circ\Phi
            =
            \Phi^F\circ(pA).
        \]
        This follows from
        \[
        \begin{aligned}
            \Phi^F\circ(pA)
            &=
            (Bp)\circ
            \Phi\circ
            (p^{-1}A)\circ
            (pA)
            \\
            &=
            (Bp)\circ\Phi.
        \end{aligned}
        \]

        The two state maps \(p\) and \(p^{-1}\) are inverse. Since both
        simulations have identity output-comparison components, their vertical
        composites also have identity output comparisons. Hence the two strict
        simulations are mutually inverse.

        Therefore every map-representable Mealy morphism is isomorphic in its
        Mealy hom-category to a canonical one arising from an internal
        functor.

        \item \textbf{Simulations correspond to internal natural
        transformations.}
        By
        Theorem~\ref{thm:stateless-mealy-simulations-natural-transformations},
        simulations
        \[
            F_{0!}\Rightarrow G_{0!}
        \]
        between canonical map-representable Mealy morphisms are exactly
        internal natural transformations
        \[
            F\Rightarrow G.
        \]
        Thus the correspondence between internal functors and canonical
        map-representable Mealy morphisms is locally an isomorphism of
        categories, and the inclusion of canonical carriers into all
        map-representable carriers is locally essentially surjective.

        \item \textbf{Compare horizontal composition.}
        Let
        \[
            F:\mathbb A\to\mathbb B,
            \qquad
            G:\mathbb B\to\mathbb C
        \]
        be internal functors. The pipeline composite of their canonical
        carriers has apex
        \[
            A_0\times_{F_0,B_0,1_{B_0}}B_0,
        \]
        which is canonically isomorphic to \(A_0\). Under this pullback
        isomorphism, the composite carrier is identified with
        \[
            (G_0F_0)_!
            =
            \left(
                A_0
                \xleftarrow{1_{A_0}}
                A_0
                \xrightarrow{G_0F_0}
                C_0
            \right).
        \]
        The pipeline transition-output structure transports to the arrow map
        \[
            G_1F_1:A_1\to C_1.
        \]
        Hence pipeline composition agrees with ordinary functor composition
        up to the canonical carrier isomorphism.

        Under the same identifications, the horizontal-composition formula
        for simulations becomes the usual horizontal pasting formula for
        internal natural transformations. The coherence isomorphisms are
        inherited from the canonical pullback associators and unitors.

        \item \textbf{Establish the biequivalence.}
        The assignment
        \[
            F\longmapsto F_{0!}
        \]
        together with the identification of internal natural transformations
        with stateless simulations is locally fully faithful on canonical
        carriers. Every map-representable 1-cell is isomorphic to one in its
        image by the strict simulations constructed above. The canonical
        pullback isomorphisms supply the pseudofunctorial composition and unit
        constraints. Therefore
        \[
            \mathbf{Cat}(\mathcal E)
            \simeq
            \mathbf{Mly}_{\mathrm{map}}(\mathcal E)
        \]
        as bicategories.

        \item \textbf{Recover the strict 2-category.}
        Restrict to canonical carriers \(F_{0!}\) and define their composition
        by transporting pipeline composition along the canonical isomorphism
        \[
            G_{0!}\circ F_{0!}
            \cong
            (G_0F_0)_!.
        \]
        Define the transported composite to be the canonical carrier of the
        ordinary composite internal functor \(GF\), and perform the analogous
        transport for simulations.

        With this transported composition, identities and associativity are
        those of ordinary internal functors and internal natural
        transformations, and hence hold strictly. The resulting strict
        2-category is the usual
        \[
            \mathbf{Cat}(\mathcal E).
        \]
    \end{enumerate}
\end{proof}

\section{Examples}
\label{sec:examples-mealy-morphisms}

Unless otherwise stated, the examples in this section are taken in \(\mathcal E=\mathbf{Set}\). Thus internal categories are small categories. Accordingly, all categories appearing below are understood to be small. In particular, when using categories such as finite sets and injections, we use a fixed small skeleton.

The first example records ordinary deterministic Mealy machines as a special case. The following examples exhibit genuinely many-object, stateful instances of the internal-category construction.

\begin{example}[Classical deterministic Mealy machines]
\label{example:classical-deterministic-mealy-machines}
    Let \(\Sigma\) be an input alphabet and \(\Gamma\) an output alphabet. Let
    \[
        \Sigma^\ast,
        \qquad
        \Gamma^\ast
    \]
    be the free monoids. Regard them as one-object categories using the one-object convention fixed earlier:
    \[
        b\circ a=ba.
    \]
    
    A deterministic Mealy machine consists of a state set \(S\), a transition function
    \[
        \tau:S\times\Sigma\to S,
    \]
    and an output function
    \[
        o:S\times\Sigma\to\Gamma.
    \]
    It extends uniquely to functions
    \[
        S\times\Sigma^\ast\to S,
        \qquad
        \Sigma^\ast\times S\to\Gamma^\ast,
    \]
    written
    \[
        s\cdot w,
        \qquad
        \omega(w,s),
    \]
    such that
    \[
        s\cdot\varepsilon=s,
        \qquad
        \omega(\varepsilon,s)=\varepsilon,
    \]
    and
    \[
        s\cdot(uv)=(s\cdot u)\cdot v,
        \qquad
        \omega(uv,s)=\omega(u,s)\omega(v,s\cdot u).
    \]
    The sequential execution diagram is
    \[
        \begin{tikzcd}[column sep=large]
                s
                    \arrow[r, "{u/\omega(u,s)}"]
            &
                s\cdot u
                    \arrow[r, "{v/\omega(v,s\cdot u)}"]
            &
                s\cdot(uv).
        \end{tikzcd}
    \]
    These are precisely the one-object Mealy laws for a Mealy morphism
    \[
        \Sigma^\ast\rightsquigarrow\Gamma^\ast.
    \]
    
    There is one minor convention to record. A general one-object Mealy morphism
    \[
        \Sigma^\ast\rightsquigarrow\Gamma^\ast
    \]
    allows the output on a generator
    \[
        a\in\Sigma
    \]
    to be an arbitrary word
    \[
        \omega(a,s)\in\Gamma^\ast.
    \]
    A classical letter-output Mealy machine is the subclass satisfying
    \[
        \omega(a,s)\in\Gamma\subseteq\Gamma^\ast
    \]
    for every \(a\in\Sigma\) and state \(s\in S\).
    
    Thus ordinary deterministic Mealy machines are recovered by passing from alphabets to their free monoids and imposing the generator-level single-letter output condition. The many-object examples below should be read as typed and stateful generalizations of this one-object recovery.
\end{example}

\begin{example}[\texorpdfstring{The base-\(r\) carry machine}{The base-r carry machine}]
\label{example:base-r-carry-machine}
    Let \(r\geq 2\). Let the input monoid be
    \[
        M=(\mathbb N,+),
    \]
    let the state set be
    \[
        P=\{0,1,\dots,r-1\},
    \]
    and define a right action by
    \[
        x\cdot n=(x+n)\bmod r.
    \]
    Let the output monoid also be
    \[
        N=(\mathbb N,+).
    \]
    Define
    \[
        \omega(n,x)=\left\lfloor\frac{x+n}{r}\right\rfloor.
    \]
    Then the additive output cocycle law becomes the carry identity
    \[
        \left\lfloor\frac{x+m+n}{r}\right\rfloor
        =
        \left\lfloor\frac{x+m}{r}\right\rfloor
        +
        \left\lfloor
        \frac{(x+m)\bmod r+n}{r}
        \right\rfloor.
    \]
    Thus a two-step addition has the form
    \[
        \begin{tikzcd}[column sep=large]
                x
                    \arrow[r, "{m/\left\lfloor (x+m)/r\right\rfloor}"]
            &
                (x+m)\bmod r
                    \arrow[r, "{n/\left\lfloor (((x+m)\bmod r)+n)/r\right\rfloor}"]
            &
                (x+m+n)\bmod r.
        \end{tikzcd}
    \]
    Therefore this data defines a Mealy morphism
    \[
        (\mathbb N,+)\rightsquigarrow(\mathbb N,+).
    \]
    The state records a residue modulo \(r\), and the output records the carry produced by adding the input.
\end{example}

\begin{example}[Typed Mealy machines from path categories]
\label{example:typed-mealy-machines-path-categories}  
    Let \(G\) and \(H\) be small directed graphs, and let
    \[
        \mathbb A=\mathbf{Path}(G),
        \qquad
        \mathbb B=\mathbf{Path}(H)
    \]
    be their free categories. The vertices of \(G\) may be read as input modes,
    and the edges of \(G\) as typed input symbols. Similarly, the vertices of
    \(H\) may be read as output modes, and the edges of \(H\) as typed output
    symbols.
    
    A Mealy morphism
    \[
        \mathbf{Path}(G)\rightsquigarrow\mathbf{Path}(H)
    \]
    assigns to each state \(x\in P\) an input vertex
    \[
        p(x)\in G_0
    \]
    and an output vertex
    \[
        q(x)\in H_0.
    \]
    If
    \[
        \alpha:u\to v
    \]
    is a path in \(G\) and \(p(x)=v\), then execution along \(\alpha\) produces
    a new state
    \[
        \delta(\alpha,x)
    \]
    over the source vertex \(u\), together with an output path
    \[
        \omega(\alpha,x):q(\delta(\alpha,x))\to q(x)
    \]
    in \(H\). The typing diagram is
    \[
    \begin{tikzcd}[column sep=huge]
            {u}
                \arrow[r, "{\alpha}"]
        &
            {v=p(x)}
    \end{tikzcd}
    \qquad\leadsto\qquad
    \begin{tikzcd}[column sep=huge]
            {q(\delta(\alpha,x))}
                \arrow[r, "{\omega(\alpha,x)}"]
        &
            {q(x)}.
    \end{tikzcd}
    \]
    
    The multiplication law says that the output path for a concatenated input
    path is obtained by concatenating the output paths produced by the
    sequential executions:
    \[
    \begin{tikzcd}[column sep=huge]
            {u}
                \arrow[r, "{\alpha}"]
        &
            {v}
                \arrow[r, "{\beta}"]
        &
            {w=p(x)}
    \end{tikzcd}
    \qquad\leadsto\qquad
    \begin{tikzcd}[column sep=huge]
            {q(\delta(\alpha,\delta(\beta,x)))}
                \arrow[r, "{\omega(\alpha,\delta(\beta,x))}"]
        &
            {q(\delta(\beta,x))}
                \arrow[r, "{\omega(\beta,x)}"]
        &
            {q(x)}.
    \end{tikzcd}
    \]
    Thus a Mealy morphism between path categories is a typed path-transducer:
    it reads input paths backwards relative to the chosen orientation and emits
    output paths in the output graph.
    
    When each of \(G\) and \(H\) has one vertex, the standard path-composition
    convention yields the opposite free monoids relative to the one-object
    convention fixed in
    Subsection~\ref{subsec:two-running-specializations}. More precisely, if
    \(\Sigma\) and \(\Gamma\) are the respective edge sets, then the one-vertex
    path categories are the one-object categories associated to
    \[
        (\Sigma^\ast)^{\mathrm{op}}
        \qquad\text{and}\qquad
        (\Gamma^\ast)^{\mathrm{op}}.
    \]
    Word reversal gives monoid isomorphisms
    \[
        \operatorname{rev}:
        \Sigma^\ast\overset{\cong}{\longrightarrow}
        (\Sigma^\ast)^{\mathrm{op}},
        \qquad
        \operatorname{rev}:
        \Gamma^\ast\overset{\cong}{\longrightarrow}
        (\Gamma^\ast)^{\mathrm{op}}.
    \]
    Hence, after transport along word reversal, the one-vertex case recovers
    the ordinary word-monoid example.
    
    In the many-object case, the construction records source and target modes
    explicitly. This is the first genuinely internal-category extension of
    ordinary Mealy machines: inputs and outputs are no longer words in a single
    alphabet, but typed paths whose endpoints matter.
    
    A Mealy simulation between two such typed path-transducers assigns to each
    state of the second transducer a state of the first with the same input
    vertex, together with a path in \(H\) comparing the two output vertices.
    The simulation law says that these comparison paths commute with the output
    paths emitted during execution.
\end{example}

\begin{example}[Context-restriction and register-allocation machines]
\label{example:context-restriction-register-allocation-machines}
    This example gives a deterministic computer-science instance whose input and output categories are genuinely many-object and whose carrier is genuinely stateful.
    
    Fix two countably infinite sets
    \[
        \mathsf{Var}
        \qquad\text{and}\qquad
        \mathsf{Reg},
    \]
    thought of as source variables and register names. Let
    \[
        \mathbf{FinSub}(\mathsf{Var})
    \]
    be the small poset category whose objects are finite subsets
    \[
        \Gamma\subseteq \mathsf{Var},
    \]
    and whose arrows are inclusions
    \[
        a:\Gamma'\hookrightarrow \Gamma.
    \]
    Similarly, let
    \[
        \mathbf{FinSub}(\mathsf{Reg})
    \]
    be the small poset category of finite subsets of registers and inclusions.
    
    Take
    \[
        \mathbb A=\mathbf{FinSub}(\mathsf{Var}),
        \qquad
        \mathbb B=\mathbf{FinSub}(\mathsf{Reg}).
    \]
    An object \(\Gamma\) of \(\mathbb A\) is read as a finite source context, for example a set of currently live source variables. An arrow
    \[
        a:\Gamma'\hookrightarrow \Gamma
    \]
    is a context inclusion. With the orientation of this chapter, processing this arrow target-to-source restricts from the larger live context \(\Gamma\) to the smaller live context \(\Gamma'\):
    \[
        \begin{tikzcd}[column sep=large]
                \Gamma'
                    \arrow[r, hook, "a"]
            &
                \Gamma.
        \end{tikzcd}
    \]
    
    Define a state to be a bijective allocation
    \[
        \rho:\Gamma\cong R
    \]
    from a finite source context \(\Gamma\subseteq\mathsf{Var}\) to a finite register context \(R\subseteq\mathsf{Reg}\). Let \(P\) be the set of triples
    \[
        (\Gamma,R,\rho)
    \]
    where
    \[
        \Gamma\subseteq\mathsf{Var},
        \qquad
        R\subseteq\mathsf{Reg}
    \]
    are finite subsets and
    \[
        \rho:\Gamma\cong R
    \]
    is a bijection. Define
    \[
        p(\Gamma,R,\rho)=\Gamma,
        \qquad
        q(\Gamma,R,\rho)=R.
    \]
    Thus the carrier span records both the source context and the currently chosen register context:
    \[
        \begin{tikzcd}
            &
                P
                    \arrow[dl, "p"']
                    \arrow[dr, "q"]
            &
            \\
                \mathbf{FinSub}(\mathsf{Var})_0
            &
            &
                \mathbf{FinSub}(\mathsf{Reg})_0.
        \end{tikzcd}
    \]
    
    Now let
    \[
        a:\Gamma'\hookrightarrow \Gamma
    \]
    be an input arrow and let
    \[
        x=(\Gamma,R,\rho)
    \]
    be a current allocation state. Define
    \[
        R':=\rho(\Gamma')\subseteq R.
    \]
    Since the output category has finite subsets of \(\mathsf{Reg}\) as literal objects, \(R'\) is an object of \(\mathbb B\), and the inclusion
    \[
        R'\hookrightarrow R
    \]
    is a literal arrow of \(\mathbb B\). Define the updated allocation by restriction:
    \[
        \rho|_{\Gamma'}:\Gamma'\cong R'.
    \]
    Thus set
    \[
        \delta(a,(\Gamma,R,\rho))
        :=
        (\Gamma',R',\rho|_{\Gamma'}),
    \]
    and define the output arrow
    \[
        \omega(a,(\Gamma,R,\rho)):R'\hookrightarrow R
    \]
    to be the inclusion of finite register subsets.
    
    The execution square is
    \[
        \begin{tikzcd}[column sep=large, row sep=large]
                \Gamma'
                    \arrow[r, hook, "{a}"]
                    \arrow[d, "{\rho|_{\Gamma'}}"', "\sim"]
            &
                \Gamma
                    \arrow[d, "{\rho}", "\sim"']
            \\
                R'
                    \arrow[r, hook, "{\omega(a,(\Gamma,R,\rho))}"']
            &
                R.
        \end{tikzcd}
    \]
    Thus execution restricts a source allocation along a context inclusion and emits the corresponding inclusion of register contexts.
    
    The unit law is immediate. Restricting along the identity inclusion
    \[
        1_\Gamma:\Gamma\hookrightarrow \Gamma
    \]
    leaves the allocation unchanged and emits the identity inclusion:
    \[
        \delta(1_\Gamma,(\Gamma,R,\rho))=(\Gamma,R,\rho),
        \qquad
        \omega(1_\Gamma,(\Gamma,R,\rho))=1_R.
    \]
    
    The multiplication law says that restriction along a composite inclusion is the same as restriction in two stages. If
    \[
        \Gamma''
            \xrightarrow{b}
        \Gamma'
            \xrightarrow{a}
        \Gamma
    \]
    are inclusions, and if
    \[
        x=(\Gamma,R,\rho),
    \]
    then
    \[
        \delta(a\circ b,x)
        =
        \delta(b,\delta(a,x)).
    \]
    Indeed, both sides are the allocation
    \[
        \rho|_{\Gamma''}:\Gamma''\cong \rho(\Gamma'').
    \]
    The output arrows also compose strictly. Put
    \[
        R'=\rho(\Gamma'),
        \qquad
        R''=\rho(\Gamma'').
    \]
    Then
    \[
        R''\subseteq R'\subseteq R,
    \]
    and
    \[
        \omega(a\circ b,x)
        =
        \omega(a,x)\circ \omega(b,\delta(a,x))
    \]
    as inclusions
    \[
        R''\hookrightarrow R.
    \]
    The two-stage square is
    \[
        \begin{tikzcd}[column sep=large, row sep=large]
                \Gamma''
                    \arrow[r, hook, "{b}"]
                    \arrow[d, "{\rho|_{\Gamma''}}"', "\sim"]
            &
                \Gamma'
                    \arrow[r, hook, "{a}"]
                    \arrow[d, "{\rho|_{\Gamma'}}"', "\sim"]
            &
                \Gamma
                    \arrow[d, "{\rho}", "\sim"']
            \\
                R''
                    \arrow[r, hook]
            &
                R'
                    \arrow[r, hook]
            &
                R.
        \end{tikzcd}
    \]
    Therefore this data defines a Mealy morphism
    \[
        \mathbf{FinSub}(\mathsf{Var})
        \rightsquigarrow
        \mathbf{FinSub}(\mathsf{Reg}).
    \]
    
    This example is not map-representable. For a fixed source context \(\Gamma\), there may be many allocations
    \[
        \rho:\Gamma\cong R
    \]
    with different choices of register subsets \(R\subseteq\mathsf{Reg}\), or different bijections to the same register subset. Hence the state is not determined by the input interface. The example models deterministic bookkeeping for environment restriction, live-variable restriction, or register-allocation restriction, without adding nondeterministic or probabilistic effects.
\end{example}

\begin{example}[Posets as internal categories]
\label{example:posets-as-internal-categories}   
    Let \(A\) and \(B\) be small posets. Every small poset may be regarded as a small category with a unique arrow
    \[
        u\to v
    \]
    precisely when
    \[
        u\leq v.
    \]
    Thus a Mealy morphism between posets \(A\) and \(B\) consists of a state set \(P\) equipped with interface maps
    \[
        p:P\to A,
        \qquad
        q:P\to B,
    \]
    together with a transition operation defined whenever
    \[
        u\leq v
        \qquad
        \text{and}
        \qquad
        p(x)=v.
    \]
    The transition produces a new state
    \[
        \delta(u\leq v,x)
    \]
    with
    \[
        p(\delta(u\leq v,x))=u,
    \]
    and an output inequality
    \[
        q(\delta(u\leq v,x))\leq q(x)
    \]
    in \(B\). The operational picture is
    \[
        \begin{tikzcd}[column sep=large]
                u 
                    \arrow[r, "\leq"] 
            & 
                v=p(x)
        \end{tikzcd}
        \qquad\leadsto\qquad
        \begin{tikzcd}[column sep=large]
                q(\delta(u\leq v,x))
                    \arrow[r, "\leq"]
            &
                q(x).
        \end{tikzcd}
    \]
    
    The unit law says that the reflexive inequality \(v\leq v\) preserves the state and emits the reflexive inequality \(q(x)\leq q(x)\). The multiplication law says that for
    \[
        u\leq v\leq w
    \]
    and \(p(x)=w\), executing along the composite inequality
    \[
        u\leq w
    \]
    is the same as first executing along \(v\leq w\) and then along \(u\leq v\). Thus a Mealy morphism between posets is a stateful refinement process that moves backwards along input refinements while emitting corresponding output refinements.
    
    A Mealy simulation
    \[
        P\Rightarrow Q
    \]
    between such Mealy morphisms consists of a function
    \[
        \theta:Q\to P
    \]
    over \(A\), together with inequalities
    \[
        q_P(\theta y)\leq q_Q(y)
    \]
    in \(B\), compatible with execution. The comparison inequality is
    \[
        \begin{tikzcd}[column sep=large]
                q_P(\theta y)
                    \arrow[r, "\leq"]
            &
                q_Q(y).
        \end{tikzcd}
    \]
    The simulation law says that the comparison inequalities commute with the output inequalities emitted by the two machines.
    
    A map-representable Mealy morphism between posets is just a monotone map. Indeed, if the carrier is
    \[
        A\xleftarrow{1_A}A\xrightarrow{F_0}B,
    \]
    then the output inequality associated to \(u\leq v\) is forced to be
    \[
        F_0u\leq F_0v.
    \]
    A Mealy simulation between map-representable carriers is a pointwise inequality
    \[
        F_0x\leq G_0x
    \]
    natural with respect to the order. This is the posetal form of an internal natural transformation.
\end{example}

\begin{example}[Action groupoids]
\label{example:action-groupoids} 
    Let a group \(G\) act on a set \(X\) on the right, and let a group \(H\)
    act on a set \(Y\) on the right. The corresponding action groupoids
    \[
        X/\!/G,
        \qquad
        Y/\!/H
    \]
    are small categories. Their arrows are
    \[
        (x,g):x\to xg,
        \qquad
        (y,h):y\to yh.
    \]
    Composition is given by
    \[
        (xg,h)\circ(x,g)=(x,gh),
    \]
    and similarly for \(Y/\!/H\).

    A Mealy morphism
    \[
        X/\!/G\rightsquigarrow Y/\!/H
    \]
    has a state set \(P\) with interface maps
    \[
        p:P\to X,
        \qquad
        q:P\to Y.
    \]
    With the orientation of this chapter, an input arrow
    \[
        (x,g):x\to xg
    \]
    is processed target-to-source: a state \(z\in P\) with
    \[
        p(z)=xg
    \]
    is updated to a state
    \[
        \delta((x,g),z)
    \]
    with
    \[
        p(\delta((x,g),z))=x.
    \]
    The output is an arrow in the action groupoid \(Y/\!/H\). Therefore it is
    determined by an element
    \[
        h((x,g),z)\in H
    \]
    such that
    \[
        q(\delta((x,g),z))\,h((x,g),z)=q(z).
    \]
    The execution diagram is
    \[
        \begin{tikzcd}[column sep=large]
            x
                \arrow[r, "{g}"]
            &
            xg=p(z)
        \end{tikzcd}
        \qquad\leadsto\qquad
        \begin{tikzcd}[column sep=large]
            q(\delta((x,g),z))
                \arrow[r, "{h((x,g),z)}"]
            &
            q(z).
        \end{tikzcd}
    \]

    A Mealy simulation between action-groupoid Mealy morphisms assigns to each
    state of the second system a state of the first system over the same input
    object, together with an \(H\)-label comparing the corresponding output
    objects. If this comparison label is \(\alpha_z\in H\), then it satisfies
    \[
        q_P(\theta z)\alpha_z=q_Q(z).
    \]
    The simulation law is the groupoid-valued naturality equation saying that
    comparison before execution and comparison after execution yield the same
    arrow in the output action groupoid.

    With the standard composition convention for right-action groupoids, the
    one-object case corresponds to the opposite group relative to the
    one-object convention
    \[
        b\circ a=ba
    \]
    fixed in Subsection~\ref{subsec:two-running-specializations}. Inversion
    gives a group isomorphism
    \[
        G\overset{\cong}{\longrightarrow}G^{\mathrm{op}},
        \qquad
        g\longmapsto g^{-1},
    \]
    and similarly for \(H\). After transport along these isomorphisms, the
    one-object case is the right-action-plus-cocycle theory of
    Subsection~\ref{subsec:one-object-mealy-morphisms}.

    Thus this example is the multi-object groupoid analogue of that
    one-object description. The state update is a contravariant action along
    input groupoid arrows, the emitted \(H\)-label is a groupoid-valued
    cocycle, and simulations compare such cocycles by coherent state-dependent
    output arrows.
\end{example}

\begin{example}[Group-valued coboundary outputs]
\label{example:group-valued-coboundary-outputs}
Let \(M\) be a monoid acting on the right on a set \(P\), and let \(G\) be a
group. Given a function
\[
    \lambda:P\to G,
\]
define
\[
    \omega(m,x)=\lambda(x)^{-1}\lambda(x\cdot m).
\]
Then
\[
\begin{aligned}
    \omega(m,x)\omega(n,x\cdot m)
    &=
    \lambda(x)^{-1}\lambda(x\cdot m)
    \lambda(x\cdot m)^{-1}\lambda((x\cdot m)\cdot n)
    \\
    &=
    \lambda(x)^{-1}\lambda((x\cdot m)\cdot n)
    \\
    &=
    \lambda(x)^{-1}\lambda(x\cdot(mn))
    \\
    &=
    \omega(mn,x).
\end{aligned}
\]
Hence
\[
    \omega(mn,x)
    =
    \omega(m,x)\omega(n,x\cdot m),
\]
so \(\omega\) is a valid output cocycle.

The cancellation in this calculation occurs entirely in the output group:
\[
    \lambda(x\cdot m)\lambda(x\cdot m)^{-1}=1_G.
\]
Thus, when \(G\) is regarded as a one-object category, all the displayed group
elements are arrows at its unique object, and the cocycle equation is an
equality of composites of such arrows.

This example shows that, for group-valued outputs, the cocycle terminology is
literal: the output law is a nonabelian cocycle equation over a right action.

The same form describes simulations in the one-object group-valued case. If
two \(G\)-valued cocycles differ by a state-dependent coboundary, then they are
connected by a Mealy simulation. For example, when \(P=Q\) and
\(\theta=1_P\), the simulation equation reads
\[
    \alpha(x)\,\omega_P(m,x)
    =
    \omega_Q(m,x)\,\alpha(x\cdot m).
\]
Equivalently,
\[
    \omega_Q(m,x)
    =
    \alpha(x)\,\omega_P(m,x)\,\alpha(x\cdot m)^{-1}.
\]

Taking
\[
    \omega_P(m,x)=1_G
\]
and
\[
    \alpha(x)=\lambda(x)^{-1},
\]
one obtains
\[
\begin{aligned}
    \omega_Q(m,x)
    &=
    \lambda(x)^{-1}
    \bigl(\lambda(x\cdot m)^{-1}\bigr)^{-1}
    \\
    &=
    \lambda(x)^{-1}\lambda(x\cdot m).
\end{aligned}
\]
Thus the coboundary cocycle
\[
    \omega_Q(m,x)=\lambda(x)^{-1}\lambda(x\cdot m)
\]
is connected to the trivial cocycle by a simulation from the trivial cocycle
to the coboundary cocycle.
\end{example}

\section{Chapter summary}
\label{sec:chapter-summary}

The constructions of the chapter may be summarized as follows:
\[
\begin{array}{c|l}
\text{structure} & \text{interpretation} \\
\hline
\text{internal category}
&
\text{monad in }\mathbf{Span}(\mathcal E)
\\
\text{Mealy morphism}
&
\text{oriented monad morphism }P A\Rightarrow B P
\\
T_B=B\circ -
&
\text{output-comparison monad}
\\
R_A=-\circ A
&
\text{input-processing monad}
\\
\text{Mealy morphism}
&
\overline R_A\text{-algebra in }\operatorname{Kl}(T_B)
\\
\text{Mealy simulation}
&
\text{oppositely directed Kleisli algebra morphism}
\\
\text{strict Mealy simulation}
&
\text{identity-output-comparison special case}
\\
\text{one-object Mealy morphism}
&
\text{internal right monoid action with output cocycle}
\\
\text{one-object Mealy simulation}
&
\text{equivariant map with coherent cocycle comparison}
\\
\text{map-representable Mealy morphism}
&
\text{internal functor, up to canonical carrier isomorphism}
\\
\text{map-representable Mealy simulation}
&
\text{internal natural transformation, up to canonical carrier isomorphism.}
\end{array}
\]

The main structural result is the bicategory
\[
    \mathbf{Mly}(\mathcal E),
\]
whose objects are internal categories, whose 1-cells are Mealy morphisms, and
whose 2-cells are Mealy simulations. The one-object specialization is internal
under the standing finite-limit assumptions: it recovers right actions of
monoid objects equipped with output cocycles. In \(\mathbf{Set}\), the free
word-monoid case recovers classical deterministic Mealy machines after imposing
the usual single-letter output condition on generators. The map-representable
specialization recovers the ordinary 2-category of internal categories,
internal functors, and internal natural transformations up to the canonical
pullback isomorphisms inherent in the span bicategory.
